\documentclass[11pt,a4paper,leqno]{amsart}%{amsproc}


\usepackage{amssymb}
\usepackage{times}
\usepackage{mathrsfs}
\usepackage{verbatim}
%\usepackage[notref]{showkeys}

\usepackage{amsrefs,doi,color}
\newcommand\half{\frac12}

%\input{colordvi}
\theoremstyle{plain}
\newtheorem{theorem}{Theorem}[section]

\theoremstyle{remark}
\newtheorem{remark}[theorem]{Remark}
\newtheorem{question}[theorem]{Question}
\newtheorem{example}[theorem]{Example}

\theoremstyle{plain}
\newtheorem{corollary}[theorem]{Corollary}
\newtheorem{lemma}[theorem]{Lemma}
\newtheorem{proposition}[theorem]{Proposition}
\newtheorem{definition}[theorem]{Definition}
\newtheorem{hypothesis}[theorem]{Hypothesis}
\newtheorem{assumption}[theorem]{Assumption}

\newtheorem{Assumption}{Assumption}[section] %%oddzielna numeracja

\renewcommand{\theAssumption}{A.\arabic{Assumption}}



\newenvironment{bew}[2]{\removelastskip\vspace{6pt}\noindent
 {\it Proof  #1.}~\rm#2}{\par\vspace{6pt}}
\newenvironment{sketch}[2]{\removelastskip\vspace{6pt}\noindent
 {\it Main ideas of the proof  #1.}~\rm#2}{\par\vspace{6pt}}


\newenvironment{del}[1]{}{\par\vspace{6pt}}
\numberwithin{equation}{section}

%f\newcommand{\del}[1]{}

\newcommand\Ce{C_{\half}^\xi(\mathbb{R},\mathrm{E})}
\newcommand\Cx{C_{\half}^\xi(\mathbb{R},\mathrm{X})}
\newcommand\Oe{\Omega(\xi,\mathrm{E})}
\newcommand\hOe{\hat{\Omega}(\xi,\mathrm{E})}
\newcommand\Ox{\Omega(\xi,\mathrm{X})}
\newcommand\Cox{C_{\half}(\mathbb{R},\mathrm{X})}

\def\imbed{\hookrightarrow}
\def\todown{\searrow}
\def\toup{\nearrow}
\newcommand{\lb}{\langle}
\newcommand{\rb}{\rangle}

\newcommand{\rf}[1]{(\ref{#1})}
\newcommand\norm[1]{|\!|#1|\!|}
\newcommand\n[1]{|\!|\!|#1|\!|\!|}
\newcommand\SetOf[2]{\bigl\{#1\,\bigm|\,#2\bigr\}}
\newcommand\bo[1]{{\bf #1}}

\newcommand{\embed}{\hookrightarrow}
\newcommand\lin{\operatorname{lin}}
\newcommand\PG{\operatorname{PG}}
\newcommand\sgn{\operatorname{sgn}}
\newcommand\grad{\operatorname{grad}}
\newcommand\diver{\operatorname{div}}
\newcommand\curl{\operatorname{curl}}
\newcommand\rot{\operatorname{rot}}
\newcommand\supp{\operatorname{supp}}
\newcommand\rank{\operatorname{rank}}
\newcommand\Res{\operatorname{Res}}
\newcommand\AG{\operatorname{AG}}
\newcommand\id{\operatorname{id}}
\newcommand\vnl{(\!(}
\newcommand\vnr{)\!)}
\newcommand\Vertt{|\!|}
\newcommand\den{:\!=}
\newcommand{\ff}{\emph{focus-focus}}
\newcommand{\U}{\mathcal{U}}
\newcommand{\M}{\mathcal{M}}

\newcommand{\eps}{\varepsilon}
\newcommand{\al}{\alpha}
\newcommand{\la}{\lambda}
\def\Nat{{\mathbb N}}
\def\Qnu{{\mathbb Q}}
\def\Rnu{{\mathbb R}}
\def\Cnu{{\mathbb C}}
\def\Rn{{\Rnu^n~}}
\def\Rt{{\Rnu^1~}}
\def\Rd{{\Rnu^d~}}
\def\Rk{{\Rnu^k~}}
\def\Rnus{\mbox{\footnotesize\bf R}}
\def\Re{\mbox{Re\,}}
\def\Im{\mbox{Im\,}}
\def\im{\mbox{im\,}}
\def\liml{\lim\limits}
\def\suml{\sum\limits}
\def\supl{\sup\limits}
\def\maxl{\max\limits}
\def\minl{\min\limits}
\def\cupl{\mathop{\cup}\limits}
\def\capl{\mathop{\cap}\limits}
\def\intl{\int\limits}
\def\prodl{\prod\limits}
\def\supl{\mathop{\sup}\limits}

\newcommand{\bmu}{\mbox{\boldmath{$\mu$}}}
\newcommand{\cit}[1]{\cite{[#1]}}

\sloppy


\begin{document}
\title{Global solutions of the random vortex filament equation}

\author{Z. Brze\'{z}niak$^\dag$, M. Gubinelli$^\ddag$ and
M. Neklyudov$^\dag$}

\address{
$^\dag$Department of Mathematics, University of York, Heslington,
UK
\\
$^\ddag$CEREMADE \& CNRS UMR 7534, Universit\'e Paris-Dauphine,  France
}
%\date{}

\thanks{This research has been supported by the ANR Project ECRU (ANR-09-BLAN-0114-01/2)}



\date{\today}

\begin{abstract}
In this article we prove the existence of a global solution for the random
vortex filament equation. Our work gives a positive answer to a question left open in recent publications: Berselli and  Gubinelli \cite{[Bers+Gub_2007]}  showed the existence of global solution for a smooth  initial condition  while  Bessaih, Gubinelli, Russo \cite{[Bess+Gub_2005]} proved the existence of a local solution for a general  initial condition.
\end{abstract}


\maketitle

In this article we prove the existence of a global solution for the following random
vortex filament equation
\begin{eqnarray}
\frac{d\gamma}{dt} &=& u^{\gamma(t)}(\gamma(t)), \;\; t\in [0,\infty) \label{eqn:RandomFilam-1}\\
\gamma(0) &=& \gamma_0,\label{eqn:RandomFilam-2}
\end{eqnarray}
where the initial condition $\gamma_0:[0,1]\to\mathbb{R}^3$ is a  geometric $\nu$-rough path (for some  $\nu\in(\frac{1}{3},1)$),  see Assumption \ref{ass:RoughPathApproximability}. Here $\gamma:[0,\infty)\to \mathcal{D}_{\gamma_0}\subset\mathcal{C}$ is some trajectory in
the subset $\mathcal{D}_{\gamma_0}$ of $\mathcal{C}$ of continuous closed curves in
$\Rnu^3$, $u^Y$, $Y\in \mathcal{D}_{\gamma_0}\subset\mathcal{C}$ is a vector field
given by
\begin{equation}
u^Y(x)=\intl_Y\nabla\phi(x-y)\times dy.\label{eqn:RandomFilam-3}
\end{equation}
where $\phi:\Rnu^3\to\Rnu$ is a smooth function which satisfies
certain assumptions (see Hypothesis \ref{hyp:PhiAssumptions}).
The exact meaning of the line integral above and set $\mathcal{D}_{\gamma_0}$ we consider
will be explained below. Equation \eqref{eqn:RandomFilam-1} appears in the fluid dynamics in
the theory of three dimensional Euler equations.
It is well known that for two dimensional Euler equations vorticity $\vec{\omega}=\curl \vec{u}$ is transported along the flow of the liquid.
The situation changes drastically in three dimensional case. Additional "stretching" term in the equation defining vorticity
leads to possibility of blow up of the vorticity.
Furthermore, a result of Beale,
Kato, and Majda \cite{[BKM-1984]} suggests that a possible
singularity of Euler equations appears when the vorticity field of the 
fluid blows up. Consequently, understanding the behaviour of vorticity of ideal fluid is
one of the most important problems in fluid dynamics.

The properties of the motion of the vorticity has been studied for the last 150 years starting from the works of
Helmholtz \cite{Helm1858} and Kelvin \cite{Kelvin1869}. It has been suggested by Kelvin  to use Biot-Savart law
\[
\vec{u}(\vec{x})=\int\frac{\vec{x}-\vec{y}}{\vert \vec{x}-\vec{y}\vert^3}\times\vec{\omega}(\vec{y})d\vec{y},\;\;\vec{x},\vec{y}\in\mathbb{R}^3
\]
where $\times$ denotes the vector product, 
combined with assumption that vorticity is supported by some smooth curve $\gamma$
\[
\omega(\vec{x},t)=\Gamma\int\limits_0^1\delta(\vec{x}-\gamma(s,t))\frac{\partial \gamma(s,t)}{\partial s}ds,x\in\mathbb{R}^3,t\geq 0
\]
and definition of the flow
\begin{equation}\label{equa21}
\left\{
\begin{array}{rcl}
\frac{d \vec{X}_t(\vec{x})}{dt}&=&\vec{u}(\vec{X}_t(\vec{x}),t),\;\;t\geq 0,\\
\vec{X}_0(\vec{x})&=&\vec{x}.
\end{array}
\right.
\end{equation}
to formally deduce the following filament equation
\begin{equation}\label{eqn:RandomFilamClassic}
\frac{\partial \gamma}{\partial t}(s,t)=-\frac{\Gamma}{4\pi}\int\limits_0^1\frac{\gamma(s,t)-\gamma(r,t)}{\vert \gamma(s,t)-\gamma(r,t)\vert^3}\times\frac{\partial\gamma(r,t)}{\partial r}dr.
\end{equation}
The assumption that vorticity is supported by some curve is coherent with numerical simulations of 3D turbulent fluids
which show that regions of large  vorticity   have a form of a
"filament", see, for instance \cite{[Bell+Markus_1992]},
\cite{[VincMeneguzzi-1991]}.

Equation \eqref{eqn:RandomFilamClassic} has singularity when $r$ is close to $s$ and the initial curve $\gamma$ is smooth.
As a consequence, the energy of the solution of this equation given by the formula
\[
E(t)=\frac{\Gamma^2}{8\pi}\int\limits_0^1\int\limits_0^1\frac{1}{\vert \gamma(s,t)-\gamma(r,t)\vert }\frac{\partial \gamma(r,t)}{\partial r}\cdot\frac{\partial \gamma(s,t)}{\partial s}dsdr,t\geq 0
\]
is infinite for any smooth curve $\gamma(\cdot,t)$.
Hence different methods have been  employed to avoid  singularity. For instance, Gallavotti \cite{Gallavotti} motivated by finiteness of the energy integral
for Brownian Motion\footnote{defined as double stochastic integral},
considers non smooth initial curves $\gamma_0$ while
Rosenhead \cite{[Rosenhead-1930]}  has suggested to use the following model
\begin{equation}\label{eqn:RandomFilamRosenhead}
\frac{\partial \gamma}{\partial t}(s,t)=-\frac{\Gamma}{4\pi}\int\limits_0^1\frac{\gamma(s,t)-\gamma(r,t)}{(\vert \gamma(s,t)-\gamma(r,t)\vert^2+\mu^2)^{3/2}}\times\frac{\partial\gamma(r,t)}{\partial r}dr.
\end{equation}
 Problem (\eqref{eqn:RandomFilam-1}-\eqref{eqn:RandomFilam-3}) has been studied  by Berselli and Bessaih~\cite{[Bers+Bess_2002]} and then by Berselli and Gubinelli~\cite{[Bers+Gub_2007]}.
It contains equation~\eqref{eqn:RandomFilamRosenhead} considered by Rosenhead as a very particular case when
\[
\phi(\vec{x})=\frac{\Gamma}{(\vert \vec{x}\vert^2+\mu^2)^{\half}},\vec{x}\in\mathbb{R}^3,\mu>0.
\]

Equation \eqref{eqn:RandomFilam-1} is in fact  a nonlinear PDE for a function $\gamma: [0,\infty)\times[0,1]\to \mathbb{R}^3$. A natural setting for the  well-posedness of the  Cauchy problem is obtained by requiring  the vector field $u$ to be well defined and Lipshitz in the space variable.
To this effect the approach followed in~\cite{[Bers+Bess_2002]} is to set up the equation as an evolution problem in the Sobolev space $H^1$ of closed curves in $\mathbb{R}^3$ with square integrable first derivative (with respect to the parameter). This approach implies that the vector field (once lifted to $H^1$) does not allow good estimates to have global existence.
This is ultimately due to the fact that in $3d$ incompressible flows vortices strech and undergo a complex dynamics and that a priori this could lead to a explosion of the $H^1$ norm.
Such a difficulty should be compared to the more stable behavior of 2d vortex points which, under incompressible flows, are simply transported along the flow lines.
Exploiting the conservation of the kinetic energy of the flow and a control of the velocity field generated by the vortex line via the associated kinetic energy Berselli and  Gubinelli \cite{[Bers+Gub_2007]}  showed the existence of global solution to  equation~\eqref{eqn:RandomFilam-1} with initial conditions in $H^1$.
After that Bessaih, Gubinelli, Russo \cite{[Bess+Gub_2005]}, partly motivated by the random filament models suggested by Gallavotti~\cite{Gallavotti} and  Chorin~\cite{Chorin}, considered the above evolution problem when the initial condition is a random closed curve which for definiteness, in the above paper, has been taken to be the sample of a Brownian loop (a Brownian motion starting at $0\in\mathbb{R}^3$ and conditioned to return to $0\in\mathbb{R}^3$ at time $1$).
In this case it is no more possible to set up the problem in the Sobolev space $H^1$ since Brownian trajectories almost surely do not belong to this space. A more serious problem is the meaning to give to the generalized Biot-Savart relation~\eqref{eqn:RandomFilam-3} since the line integral along a Brownian trajectory is a notoriously difficult object to define. Stochastic integration (\`a la It\^o or Stratonovich) does not provide a good framework to study this problem and~\cite{[Bess+Gub_2005]} identified the natural setting to have a well posed problem by considering the evolution as an equation on the space of rough paths.

Rough path theory has been introduced by T. J. Lyons in the seminal paper~\cite{[Lyons-1998]} (see also~\cites{[LyonsQian-2002],[LionsStFlour],[FrizVictoir]}) as a way to overcome certain difficulties of stochastic integration theories and have a robust analytical framework to solve stochastic differential equations and similar problems involving integration of non-regular vector-fields.
It turns out that rough paths theory and in particular the notion of \emph{controlled paths} introduced in~\cite{[Gubinelli-2004]} allows to give a natural interpretation to the Biot-Savart relation~\eqref{eqn:RandomFilam-3} and obtain a well-posed problem. Using this approach  Bessaih, Gubinelli, Russo \cite{[Bess+Gub_2005]} could obtain existence of a local solution to  equation~\eqref{eqn:RandomFilam-1} when initial condition is a closed curve of
H\"{o}lder class with exponent $\nu\in (\frac{1}{3},1]$ (suitably lifted to rough path space).

The aim of the present paper is to extend the energy method of Berselli and Gubinelli to the rough path setting to obtain  the global existence of solution of the equation~\eqref{eqn:RandomFilam-1}  when the initial condition is a (geometric) rough path, thus completing the analysis of \cite{[Bess+Gub_2005]}.

Recently there have been some progress in the study of evolution equations in the space of (controlled) rough paths. In particular Hairer~\cite{Hairer} showed how to use rough path theory to have well-posedness of a multidimensional Burgers type equation driven by additive space-time white noise. The key technical tool to obtain these results has been the observation that the non-linear term in the Burgers equation has the same structure of the Biot-Savart relation~\eqref{eqn:RandomFilam-3} and thus can be similarly handled via rough paths techniques.


\textbf{Acknowledgments.} The present paper is an extended version of a chapter from the PhD Thesis of M.~Neklyudov.

\vfill\newpage

\section{Notation}
In this section we present the controlled path framework introduced by Gubinelli~\cite{[Gubinelli-2004]}.
Let $V$ be a fixed Banach space. We introduce  the following two objects, where $S^1$ denotes the unit circle:
\begin{eqnarray*}
C_n(V)&=&\{f\in C((S^1)^n,V):\; f(t,t,\cdots,t)=0,t\in
S^1\},\\
C_*(V)&=&\cupl_{k\in\Nat}C_k(V).
\end{eqnarray*}
The operator $\delta$ is defined by 
\begin{eqnarray*}
\delta&:&C_n(V)\to C_{n+1}(V),\\
(\delta g)(t_1,\cdots,t_{n+1})&=&\suml_{i=1}^{n+1}(-1)^ig(t_1,\cdots,\hat{t_i},\cdots,t_{n+1}),\;\; n\in\Nat
\end{eqnarray*}
satisfies the following fundamental property
$$
\delta\delta=0,
$$
where $\delta\delta$ is understood as an operator from $C_n(V)$ to
$C_{n+2}(V)$. Thus $\delta$ induces a complex and we can denote
\begin{eqnarray*}
\mathcal{Z}C_k(V)&=&C_k(V)\cap\ker\delta,\\
\mathcal{B}C_k(V)&=&C_k(V)\cap\im\delta.
\end{eqnarray*}
To avoid confusion we will use notation $\delta_n$ for operator
$\delta:C_n(V)\to C_{n+1}(V)$. Furthermore, it can be noticed that
$$
\im\delta_n=\ker\delta_{n+1},
$$
i.e. $\mathcal{Z}C_{k+1}(V)=\mathcal{B}C_k(V)$.

We will  mainly consider the cases $n=1$ and $n=2$. Then the operator $\delta$
takes the following form
$$
\delta_1g(t,s)=g(t)-g(s),\delta_2h(t,u,s)=h(t,s)-h(t,u)-h(u,s),t,u,s\in
S^1.
$$
We will use special topology in spaces $C_2(V)$ and $C_3(V)$.
Let, for $\mu,\rho>0$, 
\begin{eqnarray*}
\vert f\vert_{\mu}&=&\supl_{a,b\in S^1}\frac{\vert f(a,b)\vert_{V}}{\vert a-b\vert^{\mu}},\;\;f\in
C_2(V),
\\
C_2^{\mu}(V)&=&\{f\in C_2(V):\, \vert f\vert_{\mu}<\infty\}
\\
\vert g\vert_{\rho,\mu}&=&\supl_{a,b,c\in
S^1}\frac{\vert g(a,b,c)\vert_{V}}{\vert a-b\vert^{\rho}\vert b-c\vert^{\mu}},\;\;g\in
C_3(V),
\\
C_3^{\rho,\mu}(V)&=&\{f\in C_3(V):\, \vert f\vert_{\rho,\mu}<\infty\},\;\; \vert \suml_ih_i\vert_{\mu}:=\suml_i\vert h_i\vert_{\rho_i,\mu-\rho_i}
\\
C_3^{\mu}(V)&=&\{f\in
C_3(V): \exists (h_i), \, \exists (\rho_i)\subset (0,\mu):\;  f=\suml_ih_i,\vert f\vert_{\mu}<\infty\},
\\
\mathcal{Z}C_3^{\mu}(V)&=&C_3^{\mu}(V)\cap\mathcal{Z}C_3(V),
\\
\mathcal{Z}C_3^{1+}(V)&=&\cupl_{\mu>1}\mathcal{Z}C_3^{\mu}(V),\\
C_2^{1+}(V)&=&\cupl_{\mu>1}C_2^{\mu}(V).
\end{eqnarray*}

Then following fundamental proposition has been proved in
\cite{[Gubinelli-2004]}:
\begin{proposition}
There exists an unique linear map
$\Lambda:\mathcal{Z}C_3^{1+}(V)\to C_2^{1+}(V)$ such that
$$
\delta\Lambda=\id_{\mathcal{Z}C_3^{1+}(V)}.
$$
Furthermore, for any $\mu>1$, this map is continuous from
$\mathcal{Z}C_3^{\mu}(V)$ to $C_2^{\mu}(V)$ and we have
$$
\Vert \Lambda h\Vert_{\mu}\leq
\frac{1}{2^{\mu}-2}\Vert h\Vert_{\mu},h\in\mathcal{Z}C_3^{1+}(V).
$$
\end{proposition}
Now, we define class of paths for which rough path integral will
be defined.
\begin{definition} Assume that $\nu \in (0,1)$. 
Let us fix $X\in C^{\nu}(S^1,V)$. %, $X\notin
%C^{2\nu}(S^1,\Rnu^3)$.
We say that path $Y\in C(S^1,V)$ is weakly controlled by $X$ if
there exist functions $Z\in C^{\nu}(S^1,L(V,V))$ and $R\in
C_2^{2\nu}(V)$ such that
\begin{equation}
Y(\xi)-Y(\eta) =Z(\eta)(X(\xi)-X(\eta))+R(\xi,\eta),\xi,\eta\in
S^1,\label{eqn:WeakControlDef-1}
\end{equation}
\end{definition}
Let $\mathcal{D}_X$ be the set of pairs $(Y,Z)$, where $Y\in
C(S^1,V)$ is a path weakly controlled by $X$, and $Z\in
C^{\nu}(S^1,L(V,V))$ is such that $R\in C_2^{2\nu}(V)$, where $R$
is defined by representation \eqref{eqn:WeakControlDef-1}, i.e. 
\begin{equation}
R(\xi,\eta)=Y(\xi)-Y(\eta)-Z(\eta)(X(\xi)-X(\eta)),\xi,\eta\in
S^1.\label{eqn:WeakControlDef-1'}
\end{equation}
Let us notice that $\mathcal{D}_X$ is a vector space. Let us define  semi-norm $\Vert \cdot\Vert_{\mathcal{D}_X}$  in
$\mathcal{D}_X$ by 
\begin{equation}
\Vert (Y,Z)\Vert_{\mathcal{D}_X}=\vert Z\vert_{C^{\nu}}+\Vert R\Vert_{C_2^{2\nu}},\label{eqn:D_XNorm-1}
\end{equation}
where
Furthermore, let us define a  norm $\Vert \cdot\Vert_{\mathcal{D}_X}^\ast$  in
$\mathcal{D}_X$ by 
\begin{equation}
\Vert (Y,Z)\Vert_{\mathcal{D}_X}^\ast=\Vert Y\Vert_{\mathcal{D}_X}+\vert Y\vert_{C(S^1,V)}.\label{eqn:D_XNorm-2}
\end{equation}
one can prove that
$(\mathcal{D}_X,\Vert \cdot\Vert_{\mathcal{D}_X}^\ast)$ is a Banach space.
From now on we will denote elements of $\mathcal{D}_X$ by $(Y,Y^\prime)$
and the corresponding function $R$ will be denoted by $R^Y$. We will often
omit to specify $Y^\prime$ when it is clear from the context and write
$\Vert Y\Vert_{\mathcal{D}_X}$ instead of $\Vert (Y,Y^\prime)\Vert_{\mathcal{D}_X}$.

\section{Definition and Properties of Rough Path integrals}
In this section we define rough path integral and state some of
its properties. We mainly follow \cite{[Gubinelli-2004]} and
\cite{[Bess+Gub_2005]}.  We assume that $V=\mathbb{R}^3$. 
%\begin{example}
%insert Brownian loops
%\end{example}


\begin{definition}
Let $\Pi:\mathcal{D}_X\ni (Y,Z)\mapsto Y\in C(S^1,\Rnu^3)$ be the
natural projection.
\end{definition}
We will need following properties of the $\mathcal{D}_X$, see
\cite{[Gubinelli-2004]}.
\begin{lemma}\label{lem:HolderProp}
$\Pi(\mathcal{D}_X)\subset C^{\nu}(S^1,\Rnu^3)$.
\end{lemma}
\begin{proof}[Proof of Lemma \ref{lem:HolderProp}]
It follows from equality \eqref{eqn:WeakControlDef-1}
\begin{equation}
\Vert Y\Vert_{C^{\nu}}\leq
\Vert Y\Vert_{\mathcal{D}_X}^*(1+\Vert X\Vert_{C^{\nu}}).\label{eqn:HolderIneq}
\end{equation}
The proof is complete.
\end{proof}
\begin{lemma}\label{lem:RegularMap}
Let $\phi\in C^2(\Rnu^3,\Rnu^3)$ and $(Y,Z)\in \mathcal{D}_X$.
Then
\begin{equation}
(W,W^\prime):=(\phi(Y),\phi^\prime(Y)Z)\in
\mathcal{D}_X\label{eqn:RegularMap-1}
\end{equation}
and the remainder has the following representation
\begin{eqnarray}
R^{W}(\xi,\eta)&=&\phi^\prime(Y(\xi))R(\xi,\eta)+
(Y(\eta)-Y(\xi))
\nonumber
\\&+&\intl_0^1[\nabla\phi(Y(\xi)+r(Y(\eta)-Y(\xi)))-\nabla\phi(Y(\xi))]dr,\;\;\xi,\eta\in
S^1.
\label{eqn:RegularMap-2}
\end{eqnarray}
where $R$ is the remainder for $Y$ w.r.t. $X$ given by
\eqref{eqn:WeakControlDef-1'}. Furthermore, there exists a
constant $K\geq 1$ such that
\begin{equation}
\Vert \phi(Y)\Vert_{\mathcal{D}_X}\leq
K\Vert \nabla\phi\Vert_{C^1}\Vert Y\Vert_{\mathcal{D}_X}(1+\Vert Y\Vert_{\mathcal{D}_X})(1+\Vert X\Vert_{C^{\nu}})^2.\label{eqn:RegularMap-3}
\end{equation}
Moreover, if $(\tilde{Y},\tilde{Z})\in \mathcal{D}_{\tilde{X}}$
and
$$
(\tilde{W},\tilde{W}^\prime):=(\phi(\tilde{Y}),\phi^\prime(\tilde{Y})\tilde{Z})
$$
then
\begin{equation}
\begin{split}
\label{eqn:RegularMap-4}
\vert W^\prime-\tilde{W}^\prime\vert_{C^{\nu}}+\vert R^W-R^{\tilde{W}}\vert_{C_2^{2\nu}}+\vert W-\tilde{W}\vert_{C^{\nu}}\leq\\
C(\vert X-\tilde{X}\vert_{C^{\nu}}+\vert Y^\prime-\tilde{Y}^\prime\vert_{C^{\nu}}+\vert R^Y-R^{\tilde{Y}}\vert_{C_2^{2\nu}}+\vert Y-\tilde{Y}\vert_{C^{\nu}})
\end{split}
\end{equation}
with
\begin{equation}
C=K\Vert \phi\Vert_{C^3}(1+\Vert X\Vert_{C^{\nu}}+\Vert \tilde{X}\Vert_{C^{\nu}})^3\vert (1+\Vert Y\Vert_{\mathcal{D}_X}+\Vert \tilde{Y}\Vert_{\mathcal{D}_{\tilde{X}}})^2.\label{eqn:RegularMap-4'}
\end{equation}
In the case $X=\tilde{X}$ we have
\begin{equation}
\label{eqn:RegularMap-4''}
\begin{split}
\Vert \phi(Y)-\phi(\tilde{Y})\Vert_{\mathcal{D}_X}\leq
& K\Vert \nabla\phi\Vert_{C^2}\Vert Y\Vert_{\mathcal{D}_X}\\
(1+&\Vert Y\Vert_{\mathcal{D}_X}+\Vert \tilde{Y}\Vert_{\mathcal{D}_X})^2(1+\Vert X\Vert_{C^{\nu}})^4
\Vert Y-\tilde{Y}\Vert_{\mathcal{D}_X}.
\end{split}
\end{equation}
\end{lemma}
\begin{proof}[ Proof of Lemma \ref{lem:RegularMap}]
See \cite{[Gubinelli-2004]}, Proposition $4$ for all statements of
the Lemma, except \eqref{eqn:RegularMap-2} (which is actually also
proven, though not stated explicitly). Let us show
\eqref{eqn:RegularMap-2}. Denote
$y(r)=Y(\xi)+r(Y(\eta)-Y(\xi)),r\in[0,1]$. Then
\begin{eqnarray}
\label{eqn:RegularMap-2proof}
&&
\phi(y(1))-\phi(y(0))=\intl_0^1\phi^\prime(y(r))y^\prime(r)dr
\\
\nonumber
&=&\suml_k(Y^k(\eta)-Y^k(\xi))\intl_0^1\frac{\partial\phi}{\partial
x_k}(y(r))dr\\
&=&\suml_k\frac{\partial\phi}{\partial
x_k}(Y(\xi))(Y^k(\eta)-Y^k(\xi))\nonumber
\\
&+&\suml_k(Y^k(\eta)-Y^k(\xi))\intl_0^1[\frac{\partial\phi}{\partial
x_k}(y(r))-\frac{\partial\phi}{\partial x_k}(Y(\xi))]dr\nonumber
\\
&=&\suml_{k,l}\frac{\partial\phi}{\partial
x_k}(Y(\xi))(Y^\prime)^{kl}(X^l(\eta)-X^l(\xi))
+\suml_{k}\frac{\partial\phi}{\partial
x_k}(Y(\xi))(R^Y)^k(\xi,\eta)\nonumber
\\
&+&\suml_k(Y^k(\eta)-Y^k(\xi))\intl_0^1[\frac{\partial\phi}{\partial
x_k}(y(r))-\frac{\partial\phi}{\partial
x_k}(Y(\xi))]dr,
\nonumber
\end{eqnarray}
and the result follows.
\end{proof}
Now we will define an integral of a weakly controlled by $X$ path $Y$  w.r.t.
to another weakly controlled by $X$ path $Z$. For this aim  we will need one more
definition.
\begin{definition}\label{def:RoughPath}
Assume that  $\nu>\frac{1}{3}$. We say that couple
$\mathbb{X}=(X,\mathbb{X}^2)$, $X\in C^{\nu}(S^1,\Rnu^3)$,
$\mathbb{X}^2\in C_2^{2\nu}(L(\Rnu^3,\Rnu^3))$ is a $\nu$-rough
path if the following condition is satisfied:
\begin{equation}\label{eqn:CompatCondition}
\mathbb{X}^{2}(\xi,\rho)-\mathbb{X}^{2}(\xi,\eta)-\mathbb{X}^{2}(\eta,\rho)
=(X(\xi)-X(\eta))\otimes(X(\eta)-X(\rho)), \xi,\eta,\rho\in S^1
\end{equation}
\end{definition}
\begin{remark}
If $\nu>1$ and $\mathbb{X}$ is a $\nu$-rough path, then
$\mathbb{X}$ is identically pair of constants $(X(0),0)$. Indeed,
in this case $X$ is H\"older function with exponent more than $1$
i.e. constant $X(0)$. Hence, $\mathbb{X}^{2}=0$.
\end{remark}
\begin{remark}
If $\nu\in(\frac{1}{2},1]$ then $\mathbb{X}^2$, the second
component of a $\nu$-rough path $\mathbb{X}=(X,\mathbb{X}^2)$, is
uniquely determined by its first component. Indeed, for $i,j=1,2,3$, 
\begin{equation}
\mathbb{X}^{2,ij}(\xi,\eta)=\intl_{\xi}^{\eta}(X^i_{\rho}-X^i_{\eta})dX^j_{\rho},\xi,\eta\in
S^1,\label{eqn:YoungIntegralCase}
\end{equation}
where the integral is understood in the sense of Young, see
\cite{[Young-1936]}. One can show that $\mathbb{X}^{2}$ defined by
formula \eqref{eqn:YoungIntegralCase} satisfies conditions of
Definition \ref{def:RoughPath}. Let us show the uniqueness of
$\mathbb{X}^{2}$. Assume that there exists another
$\mathbb{X}_1^{2}$ which satisfies definition \ref{def:RoughPath}.
Put $G(\xi)=\mathbb{X}^{2}(\xi,0)-\mathbb{X}_1^{2}(\xi,0)$. Then
by condition \ref{eqn:CompatCondition}
$$
\mathbb{X}^{2}(\xi,\rho)-\mathbb{X}_1^{2}(\xi,\rho)=G(\xi)-G(\rho),
$$
and, since $\mathbb{X}^{2}\in C_2^{2\nu}$, $G$ is a H\"{o}lder
function of order bigger  than $1$. Hence,  $G=0$. Therefore,
$\mathbb{X}_1^{2}=\mathbb{X}^{2}$.
\end{remark}
Note that by identity \eqref{eqn:CompatCondition} it follows that
$\mathbb{X}^2(\xi,\xi)=0,\xi\in S^1$.
\begin{assumption}\label{ass:RoughPathApproximability}
We say that our $\nu$-rough path $(X,\mathbb{X}^2)$  is an
geometric $\nu$-rough path if there exist a sequence
$(X_n,\mathbb{X}_n^2)$ such that
$$
X_n\in C^{\infty}(S^1,\Rnu^3),
$$
$$
\mathbb{X}_n^{2,ij}(\xi,\eta)=\intl_{\xi}^{\eta}(X_n^i({\rho})-X_n^i({\eta}))dX_n^j({\rho}),\xi,\eta\in
S^1,i,j=1,2,3,
$$
and
\begin{equation}
\lim_{n\to\infty}\big[\vert X_n-X\vert_{C^{\nu}}+\vert \mathbb{X}_n^2-\mathbb{X}^2\vert_{C_2^{2\nu}}\big]=
0.\label{eqn:RoughPathTopology}
\end{equation}
\end{assumption}
\begin{example}
Let $\{B_t\}_{t\in [0,1]}$ be the standard $3$-dimensional Brownian
bridge such that $B_0=B_1=x_0$ and let $\mathbb{B}^{2,ij}$, $i,j=1,2,3$,  be the
area processes defined by
$$
\mathbb{B}^{2,ij}(\xi,\eta)=\intl_{\xi}^{\eta}(B^i_{\rho}-B^i_{\eta})dB^j_{\rho},
$$
where the integral can be understood either in the Stratonovich or in the It\^o
sense. Then, the couple $(B,\mathbb{B}^2)$ is a $\nu$-rough path see \cite{[Bess+Gub_2005]}*{p.1849}. Moreover, if the integral is understood in Stratonovich sense it is
geometric $\nu$-rough path.  Indeed, it follows from Theorem
3.1 in \cite{[Friz-2005]} that one can approximate $X$ with
piecewise linear dyadic $X_n^\prime$ in the sense of assumption
\ref{ass:RoughPathApproximability}a.s..
\end{example}
From now on we suppose that the geometric $\nu$-rough path
$\mathbb{X}=(X,\mathbb{X}^2)$ and the corresponding Banach space
$\mathcal{D}_X$ are fixed.  For a  finite
partition $\pi=\{\xi_0=\xi<\xi_1<\cdots<\xi_n=\eta\}$ be   of the interval  $[\xi,\eta]$, let  $d(\pi)=\supl_i\vert \xi_{i+1}-\xi_i\vert $ denote the mesh of the partition $\pi$.
\begin{lemma}\label{lem:IntegralDef}
Let  If $Y,Z\in \mathcal{D}_X$ then the limit
\begin{equation}
\liml_{d(\pi)\to
0}\sum_{i=0}^{n-1}[Y(\xi_i)(Z(\xi_{i+1})-Z(\xi_i))+Y^\prime(\xi_i)Z^\prime(\xi_i)\mathbb{X}^2(\xi_{i+1},\xi_i)]\label{eqn:IntegralDef-1}
\end{equation}
exists and is denoted by definition by
$$
\intl_{\xi}^{\eta}YdZ.
$$
\end{lemma}
\begin{proof}[ Proof of Lemma \ref{lem:IntegralDef}]
See \cite{[Gubinelli-2004]}, Theorem $1$.
\end{proof}
\begin{remark}\label{rem:SmoothCase}
In the case of $\nu>\frac{1}{2}$ the line integral defined in the
Lemma \ref{lem:IntegralDef} is reduced to the Young definition of the
line integral $\int YdZ$. Indeed, it is enough to notice that
second term in formula \eqref{eqn:IntegralDef-1} is of the order
$O(\vert \xi_{i+1}-\xi_i\vert^{2\nu})$, $2\nu>1$. Obviously, line integral
does not depend upon $Y^\prime$, $Z^\prime$ in this case.
\end{remark}
\begin{lemma}\label{lem:IntegralProp}
Assume $Y,W\in \mathcal{D}_X$, $\tilde{Y},\tilde{W}\in
\mathcal{D}_{\tilde{X}}$. Define maps $Q,\tilde{Q}:(S^1)^2\to\Rnu$ by
the following identities
\begin{equation}
Q(\eta,\xi):=\intl_{\xi}^{\eta}YdW-Y(\xi)(W(\eta)-W(\xi))-Y^\prime(\xi)W^\prime(\xi)\mathbb{X}^2(\eta,\xi),\eta,\xi\in
S^1,\label{eqn:IntegralDef-2}
\end{equation}
\begin{equation}
\tilde{Q}(\eta,\xi):=\intl_{\xi}^{\eta}\tilde{Y}d\tilde{W}-\tilde{Y}(\xi)(\tilde{W}(\eta)-\tilde{W}(\xi))-\tilde{Y}^\prime(\xi)\tilde{W}^\prime(\xi)\tilde{\mathbb{X}}^2(\eta,\xi),\eta,\xi\in
S^1. \label{eqn:IntegralDef-4}
\end{equation}
Then $Q,\tilde{Q}\in C_2^{3\nu}$. \\
Moreover, there exists constant
$C=C(\nu)>0$ such that for all $Y,W\in \mathcal{D}_X$
\begin{equation}
\Vert Q\Vert_{C_2^{3\nu}}\leq
C(1+\Vert X\Vert_{C^{\nu}}+\Vert \mathbb{X}^2\Vert_{C_2^{2\nu}})\Vert Y\Vert_{\mathcal{D}_X}\Vert W\Vert_{\mathcal{D}_X}.\label{eqn:IntegralDef-3}
\end{equation}
Furthermore,
\begin{multline}\label{eqn:IntegralDef-5}
\Vert Q-\tilde{Q}\Vert_{C_2^{3\nu}}\leq
C(1+\Vert X\Vert_{C^{\nu}}+\Vert \mathbb{X}^2\Vert_{C_2^{2\nu}})\\
((\Vert Y\Vert_{\mathcal{D}_X}+\Vert \tilde{Y}\Vert_{\mathcal{D}_{\tilde{X}}})\eps_W+
(\Vert W\Vert_{\mathcal{D}_X}+\Vert \tilde{W}\Vert_{\mathcal{D}_{\tilde{X}}})\eps_Y+\eps_X).
\end{multline}
where
\[
\eps_Y=\vert Y^\prime-\tilde{Y}^\prime\vert_{C^{\nu}}+\vert R^Y-R^{\tilde{Y}}\vert_{C_2^{2\nu}}+\vert Y-\tilde{Y}\vert_{C^{\nu}},
\]
\[
\eps_W=\vert W^\prime-\tilde{W}^\prime\vert_{C^{\nu}}+\vert R^W-R^{\tilde{W}}\vert_{C_2^{2\nu}}+\vert W-\tilde{W}\vert_{C^{\nu}},
\]
\[
\eps_X=(\Vert Y\Vert_{\mathcal{D}_X}+\Vert \tilde{Y}\Vert_{\mathcal{D}_{\tilde{X}}})(\Vert W\Vert_{\mathcal{D}_X}+\Vert \tilde{W}\Vert_{\mathcal{D}_{\tilde{X}}})
(\vert X-\tilde{X}\vert_{C^{\nu}}+\vert \mathbb{X}^2-\mathbb{\tilde{X}}^2\vert_{C_2^{2\nu}}).
\]
\end{lemma}
\begin{proof}[ Proof of Lemma \ref{lem:IntegralProp}]
See \cite{[Gubinelli-2004]}, Theorem 1. For formula
\eqref{eqn:IntegralDef-5} see   \cite{[Gubinelli-2004]}, p.104,
formula (27).
\end{proof}
By Lemmata \ref{lem:IntegralDef} and \ref{lem:RegularMap} for any
$A\in C^2(\Rnu^3,L(\Rnu^3,\Rnu^3)),Y\in\mathcal{D}_X$ we can a
define a map  $V^Y:\Rnu^3\to\Rnu$  by invoking rough path integral as follows
\begin{equation}
V^Y(x):=\intl_{S^1}A(x-Y)dY,x\in\Rnu^3.
\end{equation}
We have following bounds on its regularity:
\begin{lemma}\label{lem:ParametricRoughPath}
Let $Y\in \mathcal{D}_X$, $\tilde{Y}\in \mathcal{D}_{\tilde{X}}$,
then there exists $C_1=C_1(\nu),C_2=C_2(\mathbb{X})$ such that for
any integer $n\geq 0$,
\begin{equation}
\Vert \nabla^nV^Y\Vert_{L^{\infty}}\leq
4C_1C_2^3\Vert \nabla^{n+1}A\Vert_{C^1}\Vert Y\Vert_{\mathcal{D}_X}^2(1+\Vert Y\Vert_{\mathcal{D}_X})\label{eqn:ParametricRoughPath-1}
\end{equation}
and
\begin{multline}\label{eqn:ParametricRoughPath-2general}
\Vert \nabla^nV^Y-\nabla^nV^{\tilde{Y}}\Vert_{L^{\infty}}\leq
C(\nu)\vert A\vert_{C^{n+3}}C_X^4(1+\Vert Y\Vert_{\mathcal{D}_X}+\Vert \tilde{Y}\Vert_{\mathcal{D}_{\tilde{X}}})^3\\
(\vert X-\tilde{X}\vert_{C^{\nu}}+\vert \mathbb{X}^2-\mathbb{\tilde{X}}^2\vert_{C_2^{2\nu}}+\vert Y^\prime-\tilde{Y}^\prime\vert_{C^{\nu}}+\vert R^Y-R^{\tilde{Y}}\vert_{C_2^{2\nu}}+\vert Y-\tilde{Y}\vert_{C^{\nu}}),
\end{multline}
where
$$
C_X=1+\vert X\vert_{C^{\nu}}+\vert \tilde{X}\vert_{C^{\nu}}+\vert \mathbb{X}^2\vert_{C_2^{2\nu}}+\vert \mathbb{\tilde{X}}^2\vert_{C_2^{2\nu}}.
$$

In the case of $X=\tilde{X}$, inequality \eqref{eqn:ParametricRoughPath-2general} can be rewritten as
\begin{equation}
\Vert \nabla^nV^Y-\nabla^nV^{\tilde{Y}}\Vert_{L^{\infty}}\leq
16C_1C_2^3\Vert \nabla^{n+1}A\Vert_{C^2}\Vert Y\Vert_{\mathcal{D}_X}(1+\Vert Y\Vert_{\mathcal{D}_X})^2\Vert Y-\tilde{Y}\Vert_{\mathcal{D}_X}^*.
\label{eqn:ParametricRoughPath-2}
\end{equation}
\end{lemma}
\begin{proof}[ Proof of Lemma \ref{lem:ParametricRoughPath}]
Inequalities \eqref{eqn:ParametricRoughPath-1} and
\eqref{eqn:ParametricRoughPath-2} were proved in
\cite{[Bess+Gub_2005]}, Lemma $7$. Now we will show
\eqref{eqn:ParametricRoughPath-2general}. It is enough to consider
the case of $n=0$. By formulae \eqref{eqn:IntegralDef-3} and
\eqref{eqn:IntegralDef-4} we have
\begin{eqnarray}
V^Y-V^{\tilde{Y}}(x)&=&A(x-Y(0))(Y(1)-Y(0))
\nonumber\\
&-&A(x-\tilde{Y}(0))(\tilde{Y}(1)-\tilde{Y}(0))\nonumber\\
&+&(A(x-Y))^\prime(0)Y^\prime(0)\mathbb{X}^2(0,1)
\nonumber\\
&-&(A(x-\tilde{Y}))^\prime(0)\tilde{Y}^\prime(0)\mathbb{\tilde{X}}^2(0,1)
\nonumber\\
&+&Q^x(0,1)-\tilde{Q}^x(0,1)\nonumber
\end{eqnarray}
where $Q^x$ and $\tilde{Q}^x$ (given by formulae
\eqref{eqn:IntegralDef-3} and \eqref{eqn:IntegralDef-4}) satisfy
inequality \eqref{eqn:IntegralDef-5} and we have identified $S^1$
with $[0,1]$. Therefore, $Y(1)=Y(0)$, $\tilde{Y}(1)=\tilde{Y}(0)$.
Hence, we have
\begin{multline}\label{eqn:PRaux-1}
\vert V^Y-V^{\tilde{Y}}\vert_{L^{\infty}}\leq
\supl_x\vert (A(x-Y))^\prime(0)Y^\prime(0)\mathbb{X}^2(0,1)-
(A(x-\tilde{Y}))^\prime(0)\tilde{Y}^\prime(0)\mathbb{\tilde{X}}^2(0,1)\vert \\
+\supl_x\vert Q^x(0,1)-\tilde{Q}^x(0,1)\vert .
\end{multline}
For the first term on the r.h.s. we have
\begin{align}\label{eqn:PRaux-2}
\vert (&A(x-Y))^\prime(0)Y^\prime(0)\mathbb{X}^2(0,1)-
(A(x-\tilde{Y}))^\prime(0)\tilde{Y}^\prime(0)\mathbb{\tilde{X}}^2(0,1)\vert \\
&\leq
\vert (dA(x-Y(0))Y^\prime(0)Y^\prime(0)-dA(x-\tilde{Y}(0))\tilde{Y}^\prime(0)\tilde{Y}^\prime(0))\mathbb{X}^2(0,1)\vert \nonumber\\
&+\vert dA(x-\tilde{Y}(0))\tilde{Y}^\prime(0)\tilde{Y}^\prime(0)(\mathbb{X}^2(0,1)-\mathbb{\tilde{X}}^2(0,1))\vert \nonumber\\
&\leq \vert \mathbb{X}^2\vert_{C_2^{2\nu}}\vert dA(x-Y(0))Y^\prime(0)Y^\prime(0)-dA(x-\tilde{Y}(0))\tilde{Y}^\prime(0)\tilde{Y}^\prime(0)\vert \nonumber\\
&+ \vert A\vert_{C^2}\vert Y^\prime\vert_{L^{\infty}}^2\vert \mathbb{X}^2-\mathbb{\tilde{X}}^2\vert_{C_2^{2\nu}}\nonumber\\
&\leq
\vert \mathbb{X}^2\vert_{C_2^{2\nu}}\vert Y^\prime\vert_{L^{\infty}}^2\vert A\vert_{C^2}\vert Y^\prime-\tilde{Y}^\prime\vert_{L^{\infty}}
+\vert \mathbb{X}^2\vert_{C_2^{2\nu}}\vert A\vert_{C^1}(\vert Y^\prime\vert_{L^{\infty}}+\vert \tilde{Y}^\prime\vert_{L^{\infty}})\vert Y^\prime-\tilde{Y}^\prime\vert_{L^{\infty}}\nonumber\\
&+\vert A\vert_{C^2}\vert Y^\prime\vert_{L^{\infty}}^2\vert \mathbb{X}^2-\mathbb{\tilde{X}}^2\vert_{C_2^{2\nu}}.\quad\nonumber
\end{align}
By \eqref{eqn:IntegralDef-5} we can estimate second term as
follows
\begin{equation}
\label{eqn:PRaux-3}
\begin{aligned}
\vert Q^x-\tilde{Q}^x\vert_{C_2^{3\nu}}&\leq
C\Big[(\Vert A(x-Y)\Vert_{\mathcal{D}_X}+\Vert A(x-\tilde{Y})\Vert_{\mathcal{D}_{\tilde{X}}})\eps_Y\nonumber\\
&+
(\Vert Y\Vert_{\mathcal{D}_X}+\Vert \tilde{Y}\Vert_{\mathcal{D}_{\tilde{X}}})\eps_A+\eps_X\Big].
\end{aligned}
\end{equation}
where
\begin{align*}
\eps_Y&=\vert Y-\tilde{Y}\vert_{C^{\nu}}+\vert Y^\prime-\tilde{Y}^\prime\vert_{C^{\nu}}+\vert R^Y-R^{\tilde{Y}}\vert_{C_2^{2\nu}},\\
\eps_A&=\vert A(x-Y)-A(x-\tilde{Y})\vert_{C^{\nu}}+\vert A(x-Y)^\prime-A(x-\tilde{Y})^\prime\vert_{C^{\nu}}\\
&+\vert R^{A(x-Y)}-R^{A(x-\tilde{Y})}\vert_{C_2^{2\nu}},\\
\eps_X&=(\Vert A(x-Y)\Vert_{\mathcal{D}_X}+\Vert A(x-\tilde{Y})\Vert_{\mathcal{D}_{\tilde{X}}})\\
&\times(\Vert Y\Vert_{\mathcal{D}_X}+\Vert \tilde{Y}\Vert_{\mathcal{D}_{\tilde{X}}})
(\vert X-\tilde{X}\vert_{C^{\nu}}+\vert \mathbb{X}^2-\mathbb{\tilde{X}}^2\vert_{C_2^{2\nu}}).
\end{align*}
By formula \eqref{eqn:RegularMap-4} we can estimate $\eps_A$ as
follows
\begin{align}
\vert \eps_A\vert &\leq
K\vert A\vert_{C^3}(1+\vert X\vert_{C^{\nu}}+\vert \tilde{X}\vert_{C^{\nu}})^3(1+\vert Y\vert_{D_X}+\vert \tilde{Y}\vert_{D_{\tilde{X}}})^2\nonumber\\
&\times(\vert X-\tilde{X}\vert_{C^{\nu}}+\vert Y-\tilde{Y}\vert_{C^{\nu}}+\vert Y^\prime-\tilde{Y}^\prime\vert_{C^{\nu}}+\vert R^Y-R^{\tilde{Y}}\vert_{C_2^{2\nu}}).\label{eqn:PRaux-4}
\end{align}
By inequality \eqref{eqn:RegularMap-3} we infer that
\begin{eqnarray}
\Vert A(x-Y)\Vert_{\mathcal{D}_X}\leq
K\vert A\vert_{C^2}\vert Y\vert_{D_X}(1+\vert Y\vert_{D_X})(1+\vert X\vert_{C^{\nu}})^2,\label{eqn:PRaux-5}
\end{eqnarray}
and similarly,
\begin{eqnarray}
\Vert A(x-\tilde{Y})\Vert_{\mathcal{D}_{\tilde{X}}}\leq
K\vert A\vert_{C^2}\vert \tilde{Y}\vert_{D_{\tilde{X}}}(1+\vert \tilde{Y}\vert_{D_{\tilde{X}}})(1+\vert \tilde{X}\vert_{C^{\nu}})^2.\label{eqn:PRaux-6}
\end{eqnarray}
Therefore, combining \eqref{eqn:PRaux-3} with \eqref{eqn:PRaux-4},
\eqref{eqn:PRaux-5} and \eqref{eqn:PRaux-6} we get
\begin{multline}\label{eqn:PRaux-7}
\vert Q^x-\tilde{Q}^x\vert_{C_2^{3\nu}}\leq C(\nu)\vert A\vert_{C^{n+3}}(1+\vert X\vert_{C^{\nu}}+\vert \tilde{X}\vert_{C^{\nu}})^4(1+\Vert Y\Vert_{\mathcal{D}_X}+\Vert \tilde{Y}\Vert_{\mathcal{D}_{\tilde{X}}})^3\\
(\vert X-\tilde{X}\vert_{C^{\nu}}+\vert \mathbb{X}^2-\mathbb{\tilde{X}}^2\vert_{C_2^{2\nu}}+\vert Y^\prime-\tilde{Y}^\prime\vert_{C^{\nu}}+\vert R^Y-R^{\tilde{Y}}\vert_{C_2^{2\nu}}+\vert Y-\tilde{Y}\vert_{C^{\nu}}).
\end{multline}
Hence, the result follows from \eqref{eqn:PRaux-2} and
\eqref{eqn:PRaux-7}.
\end{proof}
We will denote for any $Y\in\mathcal{D}_X$,
$\tilde{Y}\in\mathcal{D}_{\tilde{X}}$
$$
\vert Y-\tilde{Y}\vert_D=\vert X-\tilde{X}\vert_{C^{\nu}}+\vert \mathbb{X}^2-\mathbb{\tilde{X}}^2\vert_{C_2^{2\nu}}+\vert Y^\prime-\tilde{Y}^\prime\vert_{C^{\nu}}+\vert R^Y-R^{\tilde{Y}}\vert_{C_2^{2\nu}}+\vert Y-\tilde{Y}\vert_{C^{\nu}}.
$$

\section{Random Filaments evolution problem}
Let $\mathcal{D}_{\mathbb{X},T}=C([0,T],\mathcal{D}_X)$, where has been defined after 
\eqref{eqn:WeakControlDef-1},  be a
vector space with the usual supremum norm
\begin{equation}
\vert F\vert_{\mathcal{D}_{\mathbb{X},T}}=\supl_{t\in[0,T]}\vert F(t)\vert_{\mathcal{D}_X}^*\label{eqn:NormDef}.
\end{equation}
Obviously $\mathcal{D}_{\mathbb{X},T}$ is a Banach space. Assume
also that the function $\phi$ appeared in the formula
\eqref{eqn:RandomFilam-3} satisfies following hypothesis.
\begin{hypothesis}\label{hyp:PhiAssumptions}
\begin{trivlist}
\item[(i)]$\phi:\mathbb{R}^3\to\mathbb{R}$ is even function.
\item[(ii)] the Fourier transform of $\phi$ is real and
non-negative function:
$$
\hat{\phi}(k)\geq 0,\;\;k\in\Rnu^3
$$
\item[(iii)]
$$
\intl_{\Rnu^3}(1+\vert k\vert^2)^2\hat{\phi}(k)dk<\infty
$$
\end{trivlist}
\end{hypothesis}
\begin{example}
The function $\phi=\phi_\mu$, $\mu>0$ defined by
$$\phi(\cdot)=\frac{\Gamma}{(\vert \cdot\vert^2+\mu^2)^{\half}}$$
is smooth and satisfies Hypothesis \ref{hyp:PhiAssumptions}, see p.6
of \cite{[Bers+Gub_2007]}.
\end{example}
Then the following local existence and uniqueness Theorem for
problem \eqref{eqn:RandomFilam-1}-\eqref{eqn:RandomFilam-3} has
been proved in \cite{[Bess+Gub_2005]}, see Theorem
$3$,p.1842.
\begin{theorem}\label{thm:LocalExistUniq}
Assume $\phi\in C^6(\Rnu^3,\Rnu)$, $\nu\in (\frac{1}{3},1)$,
$\mathbb{X}=(X,\mathbb{X}^2)$ is a $\nu$-rough path, $\gamma_0\in
\mathcal{D}_X$. Then there exists a time
$T_0=T_0(\nu,\vert \phi\vert_{C^5},\mathbb{X})>0$ such that the problem
\eqref{eqn:RandomFilam-1}-\eqref{eqn:RandomFilam-3} has unique
solution in the space  $\mathcal{D}_{\gamma_0,T_0}\subset\mathcal{D}_{\mathbb{X},T_0}$.
\end{theorem}
Our aim is to prove global existence of solution of the problem
\eqref{eqn:RandomFilam-1}-\eqref{eqn:RandomFilam-3} under
assumptions of Theorem \ref{thm:LocalExistUniq} and additional
hypothesis \ref{hyp:PhiAssumptions} i.e. we shall prove
\begin{theorem}\label{thm:GlobalExistUniq}
Assume that  $\gamma_0$ is a geometric $\nu$-rough path, $\phi\in
C^6(\Rnu^3,\Rnu)$ satisfies  satisfies Hypothesis
\ref{hyp:PhiAssumptions} and  $\nu\in(\frac{1}{3},1)$. Then for every $T>0$, the problem
\eqref{eqn:RandomFilam-1}-\eqref{eqn:RandomFilam-3} has unique
solution in $\mathcal{D}_{\gamma_0,T}$.
\end{theorem}
We will need the following definition.
\begin{definition}
Let $\phi\in C^4(\Rnu^3,\Rnu)$, $\gamma\in \mathcal{D}_X$. Put
\begin{equation}
\mathcal{H}_X^{\phi}(\gamma)=\frac{1}{2}\intl_{S^1}\vec{\phi}^{\gamma}(\gamma(\xi))\cdot
d\vec{\gamma}(\xi),\label{eqn:EnergyIntegralDef-1}
\end{equation}
where
\begin{equation*}
\vec{\phi}^{\gamma}(x)=\intl_{S^1}\phi(x-\gamma(\eta))d\vec{\gamma}(\eta).\label{eqn:EnergyIntegralDef-1a}
\end{equation*}
$\mathcal{H}_X^{\phi}(\gamma)$ is called the energy of path $\gamma$. We will omit $\phi$ below.
\end{definition}

\begin{remark}\label{rem:CorrectEnergyDef}
Definition
\eqref{eqn:EnergyIntegralDef-1}-\eqref{eqn:EnergyIntegralDef-1a}
is well posed. Indeed, by Lemma \ref{lem:ParametricRoughPath}
$\vec{\psi}^{\gamma}\in C^2(\Rnu^3,\Rnu^3)$ and, therefore, it
follows by Lemma \ref{lem:RegularMap} that
$\vec{\psi}^{\gamma} \circ\gamma\in \mathcal{D}_X$.
\end{remark}
\begin{remark}
Assume that $\nu>\frac{1}{2}$ and $\gamma\in C^1(S^1,\Rnu^3)$.
Then by Remark \ref{rem:SmoothCase} the line integrals in the
definition of energy are understood in the sense of Young and
\begin{equation}
\mathcal{H}^\phi_X(\gamma)=\frac{1}{2}\intl_{S^1}\intl_{S^1}\phi(\vec{\gamma}(\xi)-\vec{\gamma}(\eta))
(\frac{d\vec{\gamma}}{d\xi}(\xi),\frac{d\vec{\gamma}}{d\eta}(\eta))d\xi
d\eta. \label{eqn:EnergyIntegralDef-2}
\end{equation}
\end{remark}
\begin{lemma}\label{lem:EnergyProperties}
Assume $\phi\in C^4(\Rnu^3,\Rnu)$. Then there exists constant
$C=C(\nu,\mathbb{X})$ such that for all $\gamma\in \mathcal{D}_X$
\begin{equation}
\vert \mathcal{H}_X(\gamma)\vert \leq
C\vert \phi\vert_{C^4}\vert \gamma\vert_{\mathcal{D}_X}^4(1+\vert \gamma\vert_{\mathcal{D}_X})^2.\label{eqn:EnergyEstimAbove}
\end{equation}
Moreover, the map $\mathcal{H}_X:\mathcal{D}_X\to\Rnu$ is locally
Lipshitz i.e. for any $R>0$ there exists $C=C(R)$ such that for
any $\gamma,\tilde{\gamma}\in\mathcal{D}_X$,
$\vert \gamma\vert_{\mathcal{D}_X}\leq R$, $\vert \tilde{\gamma}\vert_{\mathcal{D}_X}\leq R$
we have
\begin{equation}
\vert \mathcal{H}_X(\gamma)-\mathcal{H}_X(\tilde{\gamma})\vert \leq
C(R)\vert \gamma-\tilde{\gamma}\vert_{\mathcal{D}_X}^*.\label{eqn:LocLipshitz}
\end{equation}
Furthermore, for any $R>0$ there exists $C=C(R)$ such that for any
$\gamma\in\mathcal{D}_X$,
$\tilde{\gamma}\in\mathcal{D}_{\tilde{X}}$,
$$
\vert \gamma\vert_{\mathcal{D}_X}\leq R,
\vert \tilde{\gamma}\vert_{\mathcal{D}_{\tilde{X}}}\leq R,
$$
$$
C_{X,\tilde{X}}=\vert X\vert_{C^{\nu}}+\vert \tilde{X}\vert_{C^{\nu}}+\vert \mathbb{X}^2\vert_{C_2^{2\nu}}+\vert \mathbb{\tilde{X}}^2\vert_{C_2^{2\nu}}<R
$$
we have
\begin{equation}
\vert \mathcal{H}_X(\gamma)-\mathcal{H}_{\tilde{X}}(\tilde{\gamma})\vert \leq
C(R)(\vert X-\tilde{X}\vert_{C^{\nu}}+\vert \mathbb{X}^2-\mathbb{\tilde{X}}^2\vert_{C_2^{2\nu}}+\vert \gamma-\tilde{\gamma}\vert_{C^{\nu}}+\vert \gamma^\prime-\tilde{\gamma}^\prime\vert_{C^{\nu}}+\vert R^{\gamma}-R^{\tilde{\gamma}}\vert_{C_2^{2\nu}}).\label{eqn:LocLipshitz-2}
\end{equation}
\end{lemma}
\begin{proof}[ Proof of Lemma \ref{lem:EnergyProperties}]
First we will show inequality \eqref{eqn:EnergyEstimAbove}. By
representation \eqref{eqn:IntegralDef-2} we have
\begin{align}
\mathcal{H}_X(\gamma)&=\frac{1}{2}(\vec{\psi}^{\gamma}(\gamma(0))(\vec{\gamma}(1)-\vec{\gamma}(0)))\nonumber\\
&+\left[d\vec{\psi}^{\gamma}(\gamma(0))\gamma^\prime(0)\right]\gamma^\prime(0)\mathbb{X}^2(1,0)+Q(0,1)\nonumber\\
&=I+II+III,\gamma\in\mathcal{D}_X.\nonumber
\end{align}
Since $\vec{\gamma}(1)=\vec{\gamma}(0)$ we infer that $I=0$.
Concerning the second term by Lemma \ref{lem:ParametricRoughPath}
we have the following estimate
\begin{align}
\vert II\vert &\leq
\Vert \mathbb{X}^2\Vert_{C_2^{2\nu}}\Vert \nabla\vec{\psi}\Vert_{L^{\infty}}\vert \gamma^\prime\vert_{L^{\infty}}^2\nonumber\\
&\leq
C(\nu,\mathbb{X})\vert \phi\vert_{C^3}\vert \gamma\vert_{\mathcal{D}_X}^4(1+\vert \gamma\vert_{\mathcal{D}_X}).\label{eqn:EnergyEstimAbove-1}
\end{align}
For third term we infer from inequality \eqref{eqn:IntegralDef-3}
\begin{equation}
\vert III\vert \leq \Vert Q\Vert_{C_2^{3\nu}}\leq
C(\nu,\mathbb{X})\vert \vec{\psi}(\vec{\gamma})\vert_{\mathcal{D}_X}\vert \vec{\gamma}\vert_{\mathcal{D}_X}.\label{eqn:EnergyEstimAbove-2}
\end{equation}
Then by Lemmas \ref{lem:RegularMap} and
\ref{lem:ParametricRoughPath} we have
\begin{align}
\vert \vec{\psi}^{\gamma}(\vec{\gamma})\vert_{\mathcal{D}_X}&\leq
C(\nu,\mathbb{X})\vert \psi^{\gamma}\vert_{C^2}\vert \vec{\gamma}\vert_{\mathcal{D}_X}(1+\vert \vec{\gamma}\vert_{\mathcal{D}_X})\nonumber\\
&\leq C(\nu,\mathbb{X})\vert \phi\vert_{C^4}\vert \vec{\gamma}\vert_{\mathcal{D}_X}^3(1+\vert \vec{\gamma}\vert_{\mathcal{D}_X})^2\label{eqn:EnergyEstimAbove-3}
\end{align}
Combining \eqref{eqn:EnergyEstimAbove-1},
\eqref{eqn:EnergyEstimAbove-2} and \eqref{eqn:EnergyEstimAbove-3}
we get inequality \eqref{eqn:EnergyEstimAbove}.

Now we will prove inequality \eqref{eqn:LocLipshitz-2}. By formula
\eqref{eqn:IntegralDef-4} we have
\begin{align}\label{eqn:EnergyLocLip-1}
\mathcal{H}_X(\gamma)-\mathcal{H}_{\tilde{X}}(\tilde{\gamma})&=\frac{1}{2}
\Big[(\nabla\psi^{\gamma}(\gamma(0))\gamma^\prime(0)\gamma^\prime(0)-\nabla\psi^{\tilde{\gamma}}(\tilde{\gamma}(0))\tilde{\gamma}^\prime(0)\tilde{\gamma}^\prime(0))
\mathbb{X}^2(1,0)\nonumber\\
&+\nabla\psi^{\tilde{\gamma}}(\tilde{\gamma}(0))\tilde{\gamma}^\prime(0)\tilde{\gamma}^\prime(0)(\mathbb{X}^2(1,0)-\mathbb{\tilde{X}}^2(1,0))
+Q(0,1)-\tilde{Q}(0,1)\Big]\!
\nonumber\\
 &=\!I\! +\! II\! +\! III
\end{align}
The first term in \eqref{eqn:EnergyLocLip-1} can be represented as
follows
\begin{eqnarray}
I &=&
(\nabla\psi^{\gamma}(\gamma(0))\gamma^\prime(0)\gamma^\prime(0)-\nabla\psi^{\tilde{\gamma}}(\tilde{\gamma}(0))\tilde{\gamma}^\prime(0)\tilde{\gamma}^\prime(0))
\mathbb{X}^2(1,0)\nonumber\\
 &=& \left[(\nabla\psi^{\gamma}(\gamma(0))-\nabla\psi^{\tilde{\gamma}}(\tilde{\gamma}(0)))\gamma^\prime(0)\gamma^\prime(0)\right.\nonumber\\
 &+& \nabla\psi^{\tilde{\gamma}}(\tilde{\gamma}(0))(\gamma^\prime(0)-\tilde{\gamma}(0))\gamma^\prime(0)\nonumber\\
 &+& \left.\nabla\psi^{\tilde{\gamma}}(\tilde{\gamma}(0))\tilde{\gamma}(0)(\gamma^\prime(0)-\tilde{\gamma}(0))\right]\mathbb{X}^2(1,0)=A+B+C\label{eqn:EnergyLocLip-2}
\end{eqnarray}
The first term in \eqref{eqn:EnergyLocLip-2} can be estimated as
follows
\begin{eqnarray}\label{eqn:EnergyLocLip-3}
\vert A\vert &=&\vert (\nabla\psi^{\gamma}(\gamma(0))-\nabla\psi^{\tilde{\gamma}}(\tilde{\gamma}(0)))\gamma^\prime(0)\gamma^\prime(0)\mathbb{X}^2(1,0)\vert \\
&\leq&
\Vert \mathbb{X}^2\Vert_{C_2^{2\nu}}\vert \gamma\vert_{\mathcal{D}_X}^2(\vert \nabla\psi^{\gamma}
(\gamma(0))\nonumber\\
&-&\nabla\psi^{\gamma}(\tilde{\gamma}(0))\vert +\vert \nabla\psi^{\gamma}(\tilde{\gamma}(0))-\nabla\psi^{\tilde{\gamma}}(\tilde{\gamma}(0))\vert )\nonumber\\
&\leq &\Vert \mathbb{X}^2\Vert_{C_2^{2\nu}}\vert \gamma\vert_{\mathcal{D}_X}^2(\vert \psi^{\gamma}\vert_{C^2}\vert \gamma(0)
\nonumber\\
&-&\tilde{\gamma}(0)\vert
+C_X^4\vert \phi\vert_{C^4}(1+\vert \gamma\vert_{\mathcal{D}_X}+\vert \tilde{\gamma}\vert_{\mathcal{D}_{\tilde{X}}})^3\vert \gamma-\tilde{\gamma}\vert_D
\nonumber\\
&\leq&
KC_X^4\vert \phi\vert_{C^4}(1+\vert \gamma\vert_{\mathcal{D}_X}+\vert \tilde{\gamma}\vert_{\mathcal{D}_{\tilde{X}}})^3\vert \gamma-\tilde{\gamma}\vert_{D},
\nonumber
\end{eqnarray}
where second inequality follows from inequality
\eqref{eqn:ParametricRoughPath-2general} and third one from
inequality \eqref{eqn:ParametricRoughPath-1}. For second term we
have by inequality \eqref{eqn:ParametricRoughPath-1}
\begin{eqnarray}
\vert B\vert &\leq&
C\Vert \mathbb{X}^2\Vert_{C_2^{2\nu}}\vert \gamma\vert_{\mathcal{D}_X}\vert \phi\vert_{C^3}\vert \tilde{\gamma}\vert_{\mathcal{D}_X}^2
(1+\vert \tilde{\gamma}\vert_{\mathcal{D}_X})\vert \gamma-\tilde{\gamma}\vert_{D}
\nonumber
\\
&\leq& CC_X(1+\vert \gamma\vert_{\mathcal{D}_X}+\vert \tilde{\gamma}\vert_{\mathcal{D}_{\tilde{X}}})^3\vert \gamma-\tilde{\gamma}\vert_D.
\label{eqn:EnergyLocLip-4}
\end{eqnarray}
Similarly, we have for third term
\begin{eqnarray}
\vert C\vert \leq
C(\nu,\mathbb{X},\vert \gamma\vert_{\mathcal{D}_X},\vert \tilde{\gamma}\vert_{\mathcal{D}_X})\vert \gamma-\tilde{\gamma}\vert_{D}.\label{eqn:EnergyLocLip-5}
\end{eqnarray}
The term $II$ in \eqref{eqn:EnergyLocLip-1} can be estimated as
follows
\begin{align*}
\vert II\vert &\leq
\vert \nabla\psi^{\tilde{\gamma}}\vert_{L^{\infty}}\vert \tilde{\gamma}\vert_{\mathcal{D}_X}^2\vert \gamma-\tilde{\gamma}\vert_D\\
&\leq
C_X^3\vert \phi\vert_{C^3}(1+\vert \gamma\vert_{\mathcal{D}_X}+\vert \tilde{\gamma}\vert_{\mathcal{D}_{\tilde{X}}})^3\vert \gamma-\tilde{\gamma}\vert_D.
\end{align*}
Thus it remains to estimate third term of equality
\eqref{eqn:EnergyLocLip-1}. We have by inequality
\eqref{eqn:IntegralDef-5}
\begin{align*}
\vert Q(0,1)-\tilde{Q}(0,1)\vert &\leq \Vert Q-\tilde{Q}\Vert_{C_2^{3\nu}}\\
&\leq
C_X\Big[(\vert \psi^{\tilde{\gamma}}(\tilde{\gamma})\vert_{\mathcal{D}_{\tilde{X}}}+
\vert \psi^{\gamma}(\gamma)\vert_{\mathcal{D}_X})\vert \gamma-\tilde{\gamma}\vert_D\\
&+(\vert \tilde{\gamma}\vert_{\mathcal{D}_{\tilde{X}}}+\vert \tilde{\gamma}\vert_{\mathcal{D}_{\tilde{X}}})\vert \psi^{\tilde{\gamma}}(\tilde{\gamma})-\psi^{\gamma}(\gamma)\vert_D\\
&+(\vert \psi^{\tilde{\gamma}}(\tilde{\gamma})\vert_{\mathcal{D}_{\tilde{X}}}+
\vert \psi^{\gamma}(\gamma)\vert_{\mathcal{D}_X})(\vert \tilde{\gamma}\vert_{\mathcal{D}_{\tilde{X}}}+\vert \tilde{\gamma}\vert_{\mathcal{D}_{\tilde{X}}})\\
&\times(\vert X-\tilde{X}\vert_{C^{\nu}}+\vert \mathbb{X}^2-\mathbb{\tilde{X}}^2\vert_{C_2^{2\nu}})\Big]
\end{align*}
By
inequality \eqref{eqn:ParametricRoughPath-1}, the term $\vert \psi^{\tilde{\gamma}}(\tilde{\gamma})\vert_{\mathcal{D}_X}$ is
bounded by the constant
$C=C(\nu,\mathbb{X},\vert \tilde{\gamma}\vert_{\mathcal{D}_X})$. Therefore, to prove
estimate \eqref{eqn:LocLipshitz} it is enough to show that there
exists constant $C=C(\nu,\mathbb{X},R)$ such that for
$\gamma,\tilde{\gamma}\in\mathcal{D}_X$ with
$\vert \gamma\vert_{\mathcal{D}_X},\vert \tilde{\gamma}\vert_{\mathcal{D}_X}\leq R$
\begin{equation}
\vert \psi^{\tilde{\gamma}}(\tilde{\gamma})-\psi^{\gamma}(\gamma)\vert_D\leq
C\vert \tilde{\gamma}-\gamma\vert_D.\label{eqn:EnergyLocLip-6}
\end{equation}
By the triangle inequality we have
\begin{align}
\vert \psi^{\tilde{\gamma}}(\tilde{\gamma})-\psi^{\gamma}(\gamma)\vert_D
&\leq\vert \psi^{\tilde{\gamma}}(\tilde{\gamma})-\psi^{\tilde{\gamma}}(\gamma)\vert_D+
\vert \psi^{\tilde{\gamma}}(\gamma)-\psi^{\gamma}(\gamma)\vert_D\nonumber\\
&=I+II.\label{eqn:EnergyLocLip-7}
\end{align}
The first term can be estimated by using inequality
\eqref{eqn:RegularMap-4} as follows
\begin{equation}
\vert I\vert \leq
KC_X^3\vert \psi^{\tilde{\gamma}}\vert_{C^3}(1+\vert \tilde{\gamma}\vert_{\mathcal{D}_{\tilde{X}}}+
\vert \gamma\vert_{\mathcal{D}_X})^2\vert \tilde{\gamma}-\gamma\vert_D.\label{eqn:EnergyLocLip-8}
\end{equation}
By inequality \eqref{eqn:ParametricRoughPath-1} we have
\begin{equation}
\vert \psi^{\tilde{\gamma}}\vert_{C^3}\leq
C\vert \phi\vert_{C^5}\vert \tilde{\gamma}\vert_{\mathcal{D}_X}^2(1+\vert \tilde{\gamma}\vert_{\mathcal{D}_X})\label{eqn:EnergyLocLip-8'}
\end{equation}
Combining \eqref{eqn:EnergyLocLip-8} and
\eqref{eqn:EnergyLocLip-8'} we get necessary estimate for $I$. It
remains to find an estimate for term $II$. By inequalities
\eqref{eqn:RegularMap-3} and \eqref{eqn:ParametricRoughPath-2} we
have
\begin{align}
II&=\vert \psi^{\tilde{\gamma}}(\gamma)-\psi^{\gamma}(\gamma)\vert_D\nonumber\\
&\leq
(1+\vert X\vert_{\nu})\vert \psi^{\tilde{\gamma}}(\gamma)-\psi^{\gamma}(\gamma)\vert_{\mathcal{D}_X}\nonumber\\
&\leq K\vert \nabla\psi^{\tilde{\gamma}}-\nabla\psi^{\gamma}\vert_{C^1}
\vert \gamma\vert_{\mathcal{D}_X}(1+\vert \gamma\vert_{\mathcal{D}_X})(1+\vert X\vert_{\nu})^3
\nonumber\\
&\leq K\vert \phi\vert_{C^5}C_X^7(1+\vert \gamma\vert_{\mathcal{D}_X}+\vert \tilde{\gamma}\vert_{\mathcal{D}_{\tilde{X}}})^5
\vert \tilde{\gamma}-\gamma\vert_D.\label{eqn:EnergyLocLip-9}
\end{align}
Hence the inequality \eqref{eqn:LocLipshitz-2} follows. Inequality
\eqref{eqn:LocLipshitz} is a consequence of inequality
\eqref{eqn:LocLipshitz-2}.
\end{proof}
\begin{corollary}\label{cor:ContinuityEnergy}
Under assumptions of Lemma \ref{lem:EnergyProperties} and
assumption \ref{ass:RoughPathApproximability} the energy function 
$\mathcal{H}_X:\mathcal{D}_X\to\Rnu$ is  continuous. Furthermore, for any $\gamma\in\mathcal{D}_X$
$$
0\leq\mathcal{H}_X(\gamma)<\infty.
$$
\end{corollary}
\begin{proof}[Proof of Corollary \ref{cor:ContinuityEnergy}]
We only need to show that $\mathcal{H}_X(\gamma)\geq 0$, for any
$\gamma\in\mathcal{D}_X$. Other statements of the Corollary easily
follow from Lemma \ref{lem:EnergyProperties}.
%It is enough to show the following formula
Fix $n\in \mathbb{N}$.
Let $C(0)=\cup_{i=1}^{n^6}C(k_i)^n$ be a partition of the cube $C(0)$ with center $0$ and length $n^2$ on the cubes $C(k_i)^n$ of the length of $\frac{1}{n}$ with centers $k_i$ and nonintersecting interiors. Define $\hat{\phi}^n(k)=\hat{\phi}(k_i)$, $k\in C(k_i)^n$ and $0$ otherwise. Consequently, define
\[
\phi^n(x)=\suml_{i=1}^{n^6}\hat{\phi}(k_i)\int\limits_{C(k_i)^n}e^{i<k,x>}dk,x\in\mathbb{R}^3
\]
Then $\phi=\lim\limits_{n\to\infty}\phi^n$ in $C_b^4$  because of the Assumption \ref{hyp:PhiAssumptions}.
Consequently,
\begin{equation}
\mathcal{H}_X^{\phi}(\gamma)=\lim\limits_{n\to\infty}\mathcal{H}_X^{\phi_n}(\gamma).\label{aux:1}
\end{equation}
Now formula
\begin{equation}
\mathcal{H}_X^{z}(\gamma)=\frac{1}{(2\pi)^3}\intl_{\Rnu^3}\hat{z}(k)\Big\vert \intl_{S^1}e^{i(k,\gamma_s)}d\gamma_s\Big\vert^2dk,\gamma\in
\mathcal{D}_X.\label{eqn:EnergyIntegralRepr-1}
\end{equation}
is correct for $z=\phi^n$, because $\phi^n$ is a sum of a finite number of Fourier modes. Therefore,
\begin{equation}
\mathcal{H}_X^{\phi_n}(\gamma)\geq 0.\label{aux:2}
\end{equation}
Thus, the  result follows from identity \eqref{aux:1} and inequality \eqref{aux:2}.
\end{proof}
Now we are going to show that energy is a local integral of motion
for problem \eqref{eqn:RandomFilam-1}-\eqref{eqn:RandomFilam-3}.
\begin{lemma}\label{lem:EnergyMotionIntegral}
Let $\gamma\in\mathcal{D}_{\mathbb{\gamma}_0,T_0}$ be a local solution of
problem \eqref{eqn:RandomFilam-1}-\eqref{eqn:RandomFilam-3} (such a solution exists by Theorem \ref{thm:LocalExistUniq}). Then
$$
\frac{d\mathcal{H}_{\gamma_0}(\gamma(s))}{ds}=0,s\in[0,T_0).
$$
\end{lemma}
\begin{proof}[Proof of Lemma \ref{lem:EnergyMotionIntegral}]
Since $\gamma(0)=\gamma_0\in \mathcal{D}_{\gamma_0}$ is a geometric rough path
there exist sequence $\{\gamma_0^n\}_{n=1}^{\infty}\in C^{\infty}(S^1,\Rnu^3)$ such that
\[
\vert \gamma_0^n-\gamma_0\vert_{C^{\nu}}+\vert \mathbb{\gamma}^2-\mathbb{\gamma}^2\vert_{C_2^{2\nu}}\to 0,n\to\infty.
\]
Now $\gamma_0^n\in \mathcal{D}_{\gamma_0^n}$, $\gamma_0\in \mathcal{D}_{\gamma_0}$. Thus we can put
$(\gamma_0^n)^\prime=(\gamma_0)^\prime=1$ and $R^{\gamma_0^n}=R^{\gamma_0}=0$.
Hence we deduce that
$$
\vert \gamma_0^n-\gamma_0\vert_D\to 0,n\to\infty.
$$
Denote by $\gamma^n\in C([0,\infty),\mathbf{H}^1(S^1,\Rnu^3))$ the
global solution of problem
\eqref{eqn:RandomFilam-1}-\eqref{eqn:RandomFilam-3} with initial
condition $\gamma_0^n$. Existence of such solution has been proved in
Theorem $2$ of \cite{[Bers+Gub_2007]}. Then according to
\cite{[Bers+Gub_2007]} (Theorem $4$, p.1846) we have
$$
\lim_{n\to \infty} \supl_{t\in[0,T_0]}\vert \gamma^n(t)-\gamma(t)\vert_D= 0.
$$
Therefore, by the continuity of the energy functional $\mathcal{H}_{\gamma_0}$ we
have
\begin{equation}
\mathcal{H}_{\gamma_0}(\gamma(s))=\liml_{n\to\infty}\mathcal{H}_{\gamma_0^n}(\gamma^n(s)),s\in[0,T_0].\label{eqn:lemEMIaux-1}
\end{equation}
Furthermore, by Lemma $2$ of \cite{[Bers+Gub_2007]}, we
have
\begin{equation}
\mathcal{H}_{\gamma_0^n}(\gamma^n(s))=\mathcal{H}_{\gamma_0^n}(\gamma_0^n),s\in[0,T_0].\label{eqn:lemEMIaux-2}
\end{equation}
As a result, combining \eqref{eqn:lemEMIaux-1} and
\eqref{eqn:lemEMIaux-2} we get statement of the Lemma.
\end{proof}
Let us recall the definition \eqref{eqn:RandomFilam-3} of the vector field $u^Y$ generated by a $\nu$-rough path $Y$
\begin{equation}
u^Y(x)=\intl_Y\nabla\phi(x-y)\times dy,Y\in
\mathcal{D}_X.\label{eqn:RandomFilam-3'}
\end{equation}
Now we will show that if energy functional of $Y$ is bounded then
associated velocity field  is a smooth function. We have
\begin{lemma}[See Lemma $3$ in \cite{[Bers+Gub_2007]}]\label{lem:VelocitySmoothness}
For any $n\in \mathbb{Z},n\geq 0$, we have following bound
\begin{equation}
\Vert \nabla^n u^{\gamma}\Vert^2_{L^{\infty}}\leq
\frac{1}{(2\pi)^{3}}\left[\intl_{\Rnu^3}\vert \vec{k}\vert^{2(1+n)}\hat{\phi}(\vec{k})d\vec{k}\right]\mathcal{H}_X(\gamma),
\gamma\in\mathcal{D}_X,\label{eqn:VelocitySmoothness}
\end{equation}
provided that the integral
$\intl_{\Rnu^3}\vert \vec{k}\vert^{2(1+n)}\hat{\phi}(\vec{k})d\vec{k}$ is
finite and $\phi\in C^{n+4}(\Rnu^3,\Rnu^3)$.
\end{lemma}
\begin{proof}[Proof of Lemma \ref{lem:VelocitySmoothness}]
For a smooth curve $\gamma$ Lemma \ref{lem:VelocitySmoothness} has been proved in
\cite{[Bers+Gub_2007]}, see  Lemma 3. In the general case, when
$\gamma\in\mathcal{D}_X$, it is enough to notice that both sides
of inequality \eqref{eqn:VelocitySmoothness} are locally Lipshitz
and therefore, continuous w.r.t. distance
$d(Y,\tilde{Y}):=\vert Y-\tilde{Y}\vert_{\mathcal{D}_X}$, $Y\in\mathcal{D}_X$,
$\tilde{Y}\in\mathcal{D}_{\tilde{X}}$. Indeed, continuity of
$\mathcal{H}_X$ has been proven in Lemma
\ref{lem:EnergyProperties} and continuity of $\Vert \nabla^n
u^{\gamma}\Vert_{L^{\infty}}$ follows from Lemma
\ref{lem:ParametricRoughPath}.
\end{proof}
Now we are ready to prove Theorem \ref{thm:GlobalExistUniq}.
\begin{proof}[Proof of Theorem \ref{thm:GlobalExistUniq}]
According to Theorem \ref{thm:LocalExistUniq} there exists unique
local solution of problem
\eqref{eqn:RandomFilam-1}-\eqref{eqn:RandomFilam-3}. Then, we can
find $T^*>0$ and a  unique maximal local solution
$\gamma:[0,T^*)\to \mathcal{D}_{\gamma_0}$ and
\begin{equation}
\liml_{t\nearrow T^*}\Vert \gamma(t)\Vert_{\mathcal{D}_{\gamma_0}}=\infty,\label{eqn:thmGEUaux-1}
\end{equation}
see e.g. \cite{[Cartan-1971]}. Notice that we will have
\begin{equation}
\frac{d\mathcal{H}_{\gamma_0}(\gamma(s))}{ds}=0,s\in[0,T^*).\label{eqn:EnergyMotionIntegral-1}
\end{equation}
Indeed, by Theorem \ref{thm:LocalExistUniq} for any $t_0\in
[0,T^*)$ there exists unique local solution $\tilde{\gamma}$ of
problem \eqref{eqn:RandomFilam-1}, \eqref{eqn:RandomFilam-3} with
initial condition $\gamma(t_0)$ on segment $[t_0,t_0+\eps_{t_0}]$
for some $\eps_{t_0}>0$. Therefore, $\gamma=\tilde{\gamma}$ on the
segment $[t_0,t_0+\eps_{t_0}]$. Hence,
$$
\frac{d\mathcal{H}_{\gamma_0}(\gamma(s))}{ds}=0,s\in[t_0,t_0+\eps_{t_0}],t_0\in[0,T^*),
$$
and identity \eqref{eqn:EnergyMotionIntegral-1} follows. We need
to show that $T^*=\infty$. Therefore, it is enough to prove
$$
\supl_{t\in[0,T^*)}\Vert \gamma(t)\Vert_{\mathcal{D}_{\gamma_0}}<\infty.
$$
Indeed, by contradiction with \eqref{eqn:thmGEUaux-1}, the result
will follow. In the rest of the proof we show such estimate. We
recall that
\begin{equation}
\gamma(t)=\gamma_0+\intl_0^tu^{\gamma(s)}(\gamma(s))ds.\label{eqn:RandomFilam-12}
\end{equation}
Firstly we have
\begin{align}
\vert \gamma(t)\vert_{L^{\infty}}&\leq\vert \gamma_0\vert_{L^{\infty}}+\intl_0^t\vert u^{\gamma(s)}\vert_{L^{\infty}}ds\nonumber\\
&\leq\vert \gamma_0\vert_{L^{\infty}}+C\intl_0^t\mathcal{H}_{\gamma_0}(\gamma(s))ds\nonumber\\
&\leq\vert \gamma_0\vert_{L^{\infty}}+C\mathcal{H}_{\gamma_0}(\gamma_0)t,t\in[0,T^*).\label{eqn:GEUEstim-1}
\end{align}
It follows from \eqref{eqn:RandomFilam-12} that
\begin{equation}
\gamma^\prime(t)=\gamma_0^\prime+\intl_0^t\nabla
u^{\gamma(s)}(\gamma(s))\gamma^\prime(s)ds,t\in[0,T^*).\label{eqn:RandomFilam-12'}
\end{equation}
Therefore, by Lemmas \ref{lem:EnergyMotionIntegral} and
\ref{lem:VelocitySmoothness}
\begin{align}
\vert \gamma^\prime(t)\vert_{L^{\infty}}&\leq\vert \gamma^\prime_0\vert_{L^{\infty}}+\intl_0^t\vert \nabla
u^{\gamma(s)}\vert_{L^{\infty}}\vert \gamma^\prime(s)\vert_{L^{\infty}}ds\nonumber\\
&\leq\vert \gamma^\prime_0\vert_{L^{\infty}}+\intl_0^tC\mathcal{H}_{\gamma_0}^{\half}(\gamma(s))\vert \gamma^\prime(s)\vert_{L^{\infty}}ds\nonumber\\
&=\vert \gamma^\prime_0\vert_{L^{\infty}}+\intl_0^tC\mathcal{H}_{\gamma_0}^{\half}(\gamma_0)\vert \gamma^\prime(s)\vert_{L^{\infty}}ds,t\in[0,T^*).
\end{align}
Then by the Gronwall inequality we infer our second estimate
\begin{equation}
\vert \gamma^\prime(t)\vert_{L^{\infty}}\leq\vert \gamma^\prime_0\vert_{L^{\infty}}e^{C\mathcal{H}_{\gamma_0}^{\half}(\gamma_0)t},t\in[0,T^*).\label{eqn:GEUEstim-2}
\end{equation}
We will need one auxiliary estimate. We have
\begin{align}
\vert \gamma(t)\vert_{C^{\nu}}&\leq\vert \gamma_0\vert_{C^{\nu}}+\intl_0^t\vert
u^{\gamma(s)}(\gamma(s))\vert_{C^{\nu}}\vert ds\nonumber\\
&\leq\vert \gamma_0\vert_{C^{\nu}}+\intl_0^t\vert \nabla
u^{\gamma(s)}\vert_{L^{\infty}}\vert \gamma(s)\vert_{C^{\nu}}ds\nonumber\\
&\leq
\vert \gamma_0\vert_{C^{\nu}}+\intl_0^tC\mathcal{H}_{\gamma_0}^{\half}(\gamma(s))\vert \gamma(s)\vert_{C^{\nu}}ds\nonumber\\
&=\vert \gamma_0\vert_{C^{\nu}}+\intl_0^tC\mathcal{H}_{\gamma_0}^{\half}(\gamma_0)\vert \gamma(s)\vert_{C^{\nu}}ds,t\in[0,T^*).
\end{align}
Thus, by the Gronwall inequality we get
\begin{equation}
\vert \gamma(t)\vert_{C^{\nu}}\leq\vert \gamma_0\vert_{C^{\nu}}e^{C\mathcal{H}_{\gamma_0}^{\half}(\gamma_0)t},t\in[0,T^*).\label{eqn:GEUEstim-aux}
\end{equation}
Now we can estimate $C^{\nu}$ norm of $\gamma^\prime$. We have
\begin{align}
\vert \gamma^\prime(t)\vert_{C^{\nu}}&\leq\vert \gamma^\prime_0\vert_{C^{\nu}}+\intl_0^t\vert \nabla
u^{\gamma(s)}(\gamma(s))\gamma^\prime(s)\vert_{C^{\nu}}ds\nonumber\\
&\leq \vert \gamma^\prime_0\vert_{C^{\nu}}+\intl_0^t(\vert \nabla
u^{\gamma(s)}\vert_{L^{\infty}}\vert \gamma^\prime(s)\vert_{C^{\nu}}+\vert \gamma^\prime(s)\vert_{L^{\infty}}\vert \nabla
u^{\gamma(s)}(\gamma(s))\vert_{C^{\nu}})ds\nonumber\\
&\leq\vert \gamma^\prime_0\vert_{C^{\nu}}+\intl_0^t(\vert \nabla
u^{\gamma(s)}\vert_{L^{\infty}}\vert \gamma^\prime(s)\vert_{C^{\nu}}+\vert \gamma^\prime(s)\vert_{L^{\infty}}\vert \nabla^2
u^{\gamma(s)}\vert_{L^{\infty}}\vert \gamma(s)\vert_{C^{\nu}})ds\nonumber\\
&\leq\vert \gamma^\prime_0\vert_{C^{\nu}}\nonumber\\
&+\intl_0^t(C\mathcal{H}_{\gamma_0}^{\half}(\gamma_0)(\vert \gamma^\prime(s)\vert_{C^{\nu}}+\vert \gamma^\prime_0\vert_{L^{\infty}}\vert \gamma_0\vert_{C^{\nu}}e^{C\mathcal{H}_{\gamma_0}(\gamma_0)s}))ds,t\in[0,T^*),
\end{align}
where last inequality follows from Lemmas
\ref{lem:EnergyMotionIntegral} and \ref{lem:VelocitySmoothness}.
Then by the Gronwall inequality we get the third estimate
\begin{equation}
\vert \gamma^\prime(t)\vert_{C^{\nu}}\leq(\vert \gamma^\prime_0\vert_{C^{\nu}}+\vert \gamma^\prime_0\vert_{L^{\infty}}\vert \gamma_0\vert_{C^{\nu}})e^{C\mathcal{H}_{\gamma_0}(\gamma_0)t},t\in[0,T^*).\label{eqn:GEUEstim-3}
\end{equation}
%and by inequality \eqref{eqn:GEUEstim-aux} we get third estimate
%\begin{equation}
%\vert \gamma^\prime(t)\vert_{C^{\nu}}\leq\vert \gamma^\prime_0\vert_{C^{\nu}}e^{C\mathcal{H}_X^{\half}(\gamma_0)te^{C\mathcal{H}_X^{\half}(\gamma_0)t}},t\in[0,T^*).\label{eqn:GEUEstim-3}
%\end{equation}
It remains to find an estimate for $\vert R^{\gamma(t)}\vert_{2\nu}$. We
have
\begin{equation}
R^{\gamma(t)}=R^{\gamma_0}+\intl_0^tR^{u^{\gamma(s)}(\gamma(s))}ds,t\in[0,T^*).
\end{equation}
By identity \eqref{eqn:RegularMap-2} we have for $s\in[0,T^*)$
\begin{align}
R^{u^{\gamma(s)}(\gamma(s))}(\xi,\eta)&=\nabla
u^{\gamma(s)}(\gamma(s,\xi))R^{\gamma(s)}(\xi,\eta)+
\suml_k(\gamma^k(s,\eta)-\gamma^k(s,\xi))\times\nonumber\\
&\intl_0^1\left[\frac{\partial u^{\gamma(s)}}{\partial
x_k}(\gamma(s,\xi)+r(\gamma(s,\eta)-\gamma(s,\xi)))-\frac{\partial
u^{\gamma(s)}}{\partial
x_k}(\gamma(s,\xi))\right]dr.\label{eqn:GEUEstim-aux3}
\end{align}
Therefore,
\begin{equation}
\vert R^{u^{\gamma(s)}(\gamma(s))}\vert_{\tilde{C}^{2\nu}}\leq \vert \nabla
u^{\gamma(s)}\vert_{L^{\infty}}\vert R^{\gamma(s)}\vert_{\tilde{C}^{2\nu}}+\frac{1}{2}\vert \gamma(s)\vert_{C^{\nu}}^2\vert \nabla^2
u^{\gamma(s)}\vert_{L^{\infty}},s\in[0,T^*).\label{eqn:GEUEstim-aux4}
\end{equation}
Thus, by inequalities \eqref{eqn:GEUEstim-aux4} and
\eqref{eqn:GEUEstim-aux} we have for $t\in[0,T^*)$
\begin{align}
\vert R^{\gamma(t)}\vert_{\tilde{C}^{2\nu}}&\leq\vert R^{\gamma_0}\vert_{\tilde{C}^{2\nu}}+\intl_0^t(\vert \nabla
u^{\gamma(s)}\vert_{L^{\infty}}\vert R^{\gamma(s)}\vert_{\tilde{C}^{2\nu}}+\frac{1}{2}\vert \gamma(s)\vert_{C^{\nu}}^2\vert \nabla^2
u^{\gamma(s)}\vert_{L^{\infty}})ds\nonumber\\
&\leq\vert R^{\gamma_0}\vert_{\tilde{C}^{2\nu}}+\intl_0^t(\vert \nabla
u^{\gamma(s)}\vert_{L^{\infty}}\vert R^{\gamma(s)}\vert_{\tilde{C}^{2\nu}}+\vert \gamma_0\vert_{C^{\nu}}e^{C\mathcal{H}_{\gamma_0}^{\half}(\gamma_0)t}\vert \nabla^2
u^{\gamma(s)}\vert_{L^{\infty}})ds\nonumber\\
&\leq\vert R^{\gamma_0}\vert_{\tilde{C}^{2\nu}}+C(\vert \gamma_0\vert_{C^{\nu}},\mathcal{H}_{\gamma_0}(\gamma_0))e^{C\mathcal{H}_{\gamma_0}^{\half}(\gamma_0)t}\nonumber\\
&+\intl_0^tC\mathcal{H}_{\gamma_0}^{\half}(\gamma_0)\vert R^{\gamma(s)}\vert_{\tilde{C}^{2\nu}}ds,\label{eqn:GEUEstim-aux5}
\end{align}
where in the last inequality we used Lemmas
\ref{lem:EnergyMotionIntegral} and \ref{lem:VelocitySmoothness}.
Hence, by the Gronwall Lemma we get
\begin{equation}
\vert R^{\gamma(t)}\vert_{\tilde{C}^{2\nu}}\leq(\vert R^{\gamma_0}\vert_{\tilde{C}^{2\nu}}+C(\vert \gamma_0\vert_{C^{\nu}},\mathcal{H}_{\gamma_0}(\gamma_0))e^{C\mathcal{H}_{\gamma_0}^{\half}(\gamma_0)t}))
e^{C\mathcal{H}_{\gamma_0}^{\half}(\gamma_0)t},t\in[0,T^*),\label{eqn:GEUEstim-4}
\end{equation}
and combining estimates \eqref{eqn:GEUEstim-1},
\eqref{eqn:GEUEstim-2}, \eqref{eqn:GEUEstim-3}, and
\eqref{eqn:GEUEstim-4} we prove following a'priori estimate
\begin{equation}
\vert \gamma(t)\vert_{\mathcal{D}_{\gamma_0}}\leq
K(1+\mathcal{H}_{\gamma_0}(\gamma_0))(1+\vert \gamma_0\vert_{\mathcal{D}_{\gamma_0}})\vert \gamma_0\vert_{\mathcal{D}_{\gamma_0}}e^{C\mathcal{H}_{\gamma_0}(\gamma_0)t},t\in[0,T^*),\label{eqn:GEUEstim-5}
\end{equation}
and the result follows.
\end{proof}

\section{Future directions of research}

It would be interesting to consider problem
\eqref{eqn:RandomFilam-1}-\eqref{eqn:RandomFilam-3} with added
white noise i.e. to consider problem
\begin{eqnarray}
d\gamma(t) &=& u^{\gamma(t)}(\gamma(t))dt+\sqrt{2\nu}dw_t,\nu>0,t\in [0,T] \label{eqn:FutureRandomFilam-1}\\
\gamma(0) &=& \gamma_0,\label{eqn:FutureRandomFilam-2}
\end{eqnarray}
where $\gamma_0$ is a geometric $\nu$-rough path, vector field of velocity $u^Y$ is given by
\eqref{eqn:RandomFilam-2} and $w_t$ is $\mathcal{D}_{\gamma_0}$--valued
Wiener process. This model would correspond to Navier-Stokes
equations rather than Euler equations. There are two possible
mathematical frameworks for the model.

First one is to make change of variables
$$
\alpha(t)=\gamma(t)-\sqrt{2\nu}w_t, t\in[0,T].
$$
Then we can fix $\{w_t\}_{t\geq 0}$ and system
\eqref{eqn:FutureRandomFilam-1}-\eqref{eqn:FutureRandomFilam-2} is
reformulated as follows
\begin{eqnarray}
\frac{d\alpha}{dt} &=& u^{\alpha(t)+\sqrt{2\nu}w_t}(\alpha(t)+\sqrt{2\nu}w_t),\nu>0,t\in [0,T] \label{eqn:FutureRandomFilam-1a}\\
\alpha(0) &=& \gamma_0.\label{eqn:FutureRandomFilam-2a}
\end{eqnarray}
Now the problem
\eqref{eqn:FutureRandomFilam-1a}-\eqref{eqn:FutureRandomFilam-2a}
is ordinary differential equation (ODE) with random coefficients
in $\mathcal{D}_{\gamma_0}$ and it can be studied by methods of the theory
of random dynamical systems, see \cite{[Arnold-98]} and
\cite{[DuanLuSchmalfus-2003]}. This approach works only in the
case of additive noise.

Second approach is to consider problem
\eqref{eqn:FutureRandomFilam-1}-\eqref{eqn:FutureRandomFilam-2} as
SDE in Banach space $\mathcal{D}_{\gamma_0}$. Then, we can consider more
general system with multiplicative noise:
\begin{eqnarray}
d\gamma(t) &=& u^{\gamma(t)}(\gamma(t))dt+\sqrt{2\nu}G(\gamma)dw_t,\nu>0,t\in [0,T] \label{eqn:FutureRandomFilam-1b}\\
\gamma(0) &=& \gamma_0.\label{eqn:FutureRandomFilam-2b}
\end{eqnarray}
The problem which appear here is to define Stochastic integral in
the Banach space $\mathcal{D}_{\gamma_0}$. Stochastic calculus in M-type 2
Banach spaces developed in works
\cite{[Dettweiler-85]}-\cite{[Dettweiler-91]}, \cite{[Brz-95]},
\cite{[Brz-03]} does not work in this situation. It seems that it
is necessary to try to alter definition of $\mathcal{D}_{\gamma_0}$ to be
able to apply the theory.

Other possible direction of research is the theory of connections
on infinite dimensional manifolds, see \cite{[Gawedzki-1988]},
\cite{[Brylinski-1992]}, \cite{[Leandre-1999]}. In
\cite{[Brz-Leandre-2000]} the authours claimed, see p. 251
therein, that it is possible to define the topological space of
Gaw{\c e}dzki's \cite{[Gawedzki-1988]} line bundle over the set of
rough loops in the sense of Lyons \cite{[Lyons-1998]}. Since the
trajectories of the Brownian loop are almost surely rough paths,
this allows us to define the topological space of Gaw{\c e}dzki's
line bundle over the Brownian bridge, because it is possible to
define the integral of a one-form over a rough path. It would be
interesting to write down a complete proof of this claim. The
theory presented in this article could be seen as a first step in
realizing such a programme.


\begin{bibdiv}
\begin{biblist}

\bib{[Arnold-98]}{book}{
   author={Arnold, Ludwig},
   title={Random dynamical systems},
   series={Springer Monographs in Mathematics},
   publisher={Springer-Verlag},
   place={Berlin},
   date={1998},
   pages={xvi+586},
   isbn={3-540-63758-3},
   review={\MR{1723992 (2000m:37087)}},
}

\bib{[BKM-1984]}{article}{
   author={Beale, J. T.},
   author={Kato, T.},
   author={Majda, A.},
   title={Remarks on the breakdown of smooth solutions for the $3$-D Euler
   equations},
   journal={Comm. Math. Phys.},
   volume={94},
   date={1984},
   number={1},
   pages={61--66},
   issn={0010-3616},
   review={\MR{763762 (85j:35154)}},
}

\bib{[Bell+Markus_1992]}{article}{
   author={Bell, John B.},
   author={Marcus, Daniel L.},
   title={Vorticity intensification and transition to turbulence in the
   three-dimensional Euler equations},
   journal={Comm. Math. Phys.},
   volume={147},
   date={1992},
   number={2},
   pages={371--394},
   issn={0010-3616},
   review={\MR{1174419 (93c:76048)}},
}

\bib{[Bers+Bess_2002]}{article}{
   author={Berselli, Luigi C.},
   author={Bessaih, Hakima},
   title={Some results for the line vortex equation},
   journal={Nonlinearity},
   volume={15},
   date={2002},
   number={6},
   pages={1729--1746},
   issn={0951-7715},
   review={\MR{1938468 (2003m:76028)}},
   doi={10.1088/0951-7715/15/6/301},
}

\bib{[Bers+Gub_2007]}{article}{
   author={Berselli, Luigi C.},
   author={Gubinelli, Massimiliano},
   title={On the global evolution of vortex filaments, blobs, and small
   loops in 3D ideal flows},
   journal={Comm. Math. Phys.},
   volume={269},
   date={2007},
   number={3},
   pages={693--713},
   issn={0010-3616},
   review={\MR{2276358 (2007m:76027)}},
   doi={10.1007/s00220-006-0142-x},
}

\bib{[Bess+Gub_2005]}{article}{
   author={Bessaih, Hakima},
   author={Gubinelli, Massimiliano},
   author={Russo, Francesco},
   title={The evolution of a random vortex filament},
   journal={Ann. Probab.},
   volume={33},
   date={2005},
   number={5},
   pages={1825--1855},
   issn={0091-1798},
   review={\MR{2165581 (2006i:60069)}},
   doi={10.1214/009117905000000323},
}

\bib{[Brylinski-1992]}{book}{
   author={Brylinski, Jean-Luc},
   title={Loop spaces, characteristic classes and geometric quantization},
   series={Progress in Mathematics},
   volume={107},
   publisher={Birkh\"auser Boston Inc.},
   place={Boston, MA},
   date={1993},
   pages={xvi+300},
   isbn={0-8176-3644-7},
   review={\MR{1197353 (94b:57030)}},
}

\bib{[Brz-95]}{article}{
   author={Brze{\'z}niak, Zdzis{\l}aw},
   title={Stochastic partial differential equations in M-type $2$ Banach
   spaces},
   journal={Potential Anal.},
   volume={4},
   date={1995},
   number={1},
   pages={1--45},
   issn={0926-2601},
   review={\MR{1313905 (95m:35213)}},
   doi={10.1007/BF01048965},
}

\bib{[Brz-Leandre-2000]}{article}{
   author={Brze{\'z}niak, Zdzis{\l}aw},
   author={L{\'e}andr{\'e}, Remi},
   title={Horizontal lift of an infinite dimensional diffusion},
   journal={Potential Anal.},
   volume={12},
   date={2000},
   number={3},
   pages={249--280},
   issn={0926-2601},
   review={\MR{1752854 (2001d:58043)}},
   doi={10.1023/A:1008622712051},
}

\bib{[Brz-03]}{article}{
   author={Brze{\'z}niak, Zdzis{\l}aw},
   title={Some remarks on It{\^o} and Stratonovich integration in 2-smooth
   Banach spaces},
   conference={
      title={Probabilistic methods in fluids},
   },
   book={
      publisher={World Sci. Publ., River Edge, NJ},
   },
   date={2003},
   pages={48--69},
   volume={XX},
   review={\MR{2083364 (2005g:60085)}},
   doi={10.1142/9789812703989\_0004},
}

\bib{[Cartan-1971]}{book}{
   author={Cartan, Henri},
   title={Differential calculus},
   note={Exercises by C. Buttin, F. Rideau and J. L. Verley;
   Translated from the French},
   publisher={Hermann},
   place={Paris},
   date={1971},
   pages={160},
}

\bib{Chorin}{book}{
   author={Chorin, Alexandre J.},
   title={Vorticity and turbulence},
   series={Applied Mathematical Sciences},
   volume={103},
   publisher={Springer-Verlag},
   place={New York},
   date={1994},
   pages={viii+174},
   isbn={0-387-94197-5},
   review={\MR{1281384 (95m:76043)}},
}

\bib{[Dettweiler-85]}{article}{
   author={Dettweiler, Egbert},
   title={ Stochastic integration of Banach space valued functions},
   book={
 	title = {Stochastic Space-Time Models and Limit Theorems},
	isbn = {9789027720382},
	language = {en},
	publisher = {Springer},
	editor = {Arnold, L.},
	editor = {Kotelenez, P.},
	year = {1985},
   },
   pages={33--79},
}


\bib{[Dettweiler-91]}{article}{
   author={Dettweiler, Egbert},
   title={Stochastic integration relative to Brownian motion on a general
   Banach space},
   language={English, with Turkish summary},
   journal={Do\u ga Mat.},
   volume={15},
   date={1991},
   number={2},
   pages={58--97},
   issn={1010-7622},
   review={\MR{1115509 (93b:60112)}},
}

\bib{[DuanLuSchmalfus-2003]}{article}{
   author={Duan, Jinqiao},
   author={Lu, Kening},
   author={Schmalfuss, Bj{\"o}rn},
   title={Invariant manifolds for stochastic partial differential equations},
   journal={Ann. Probab.},
   volume={31},
   date={2003},
   number={4},
   pages={2109--2135},
   issn={0091-1798},
   review={\MR{2016614 (2004m:60136)}},
   doi={10.1214/aop/1068646380},
}

\bib{[Friz-2005]}{article}{
   author={Friz, Peter K.},
   title={Continuity of the It\^o-map for H\"older rough paths with
   applications to the support theorem in H\"older norm},
   conference={
      title={Probability and partial differential equations in modern
      applied mathematics},
   },
   book={
      series={IMA Vol. Math. Appl.},
      volume={140},
      publisher={Springer},
      place={New York},
   },
   date={2005},
      volume={XX},
   pages={117--135},
   review={\MR{2202036 (2007f:60070)}},
   doi={10.1007/978-0-387-29371-4\_8},
}

\bib{[FrizVictoir]}{book}{
   author={Friz, Peter K.},
   author={Victoir, Nicolas B.},
   title={Multidimensional stochastic processes as rough paths},
   series={Cambridge Studies in Advanced Mathematics},
   volume={120},
   note={Theory and applications},
   publisher={Cambridge University Press},
   place={Cambridge},
   date={2010},
   pages={xiv+656},
   isbn={978-0-521-87607-0},
   review={\MR{2604669 (2012e:60001)}},
}

\bib{Gallavotti}{book}{
   author={Gallavotti, Giovanni},
   title={Foundations of fluid dynamics},
   series={Texts and Monographs in Physics},
   note={Translated from the Italian},
   publisher={Springer-Verlag},
   place={Berlin},
   date={2002},
   pages={xviii+513},
   isbn={3-540-41415-0},
   review={\MR{1872661 (2003e:76002)}},
}

\bib{[Gawedzki-1988]}{article}{
   author={Gaw{\polhk{e}}dzki, K.},
   title={Topological actions in two-dimensional quantum field theories},
   conference={
      title={Nonperturbative quantum field theory},
      address={Carg\`ese},
      date={1987},
   },
   book={
      series={NATO Adv. Sci. Inst. Ser. B Phys.},
      volume={185},
      publisher={Plenum},
      place={New York},
   },
   date={1988},
   pages={101--141},
   review={\MR{1008277 (90i:81122)}},
}

\bib{[Gubinelli-2004]}{article}{
   author={Gubinelli, M.},
   title={Controlling rough paths},
   journal={J. Funct. Anal.},
   volume={216},
   date={2004},
   number={1},
   pages={86--140},
   issn={0022-1236},
   review={\MR{2091358 (2005k:60169)}},
   doi={10.1016/j.jfa.2004.01.002},
}

\bib{Hairer}{article}{
   author={Hairer, M.},
   title={Rough stochastic PDEs},
   journal={Comm. Pure Appl. Math.},
   volume={64},
   date={2011},
   number={11},
   pages={1547--1585},
   issn={0010-3640},
   review={\MR{2832168}},
   doi={10.1002/cpa.20383},
}


\bib{Helm1858}{article}{
	title = {\"Uber Integrale der hydrodynamischen Gleichungen, welche den Wirbelbewegungen entsprechen.},
	volume = {1858},
	issn = {0075-4102, 1435-5345},
	doi = {10.1515/crll.1858.55.25},
	number = {55},
	journal = {Journal f\"ur die reine und angewandte Mathematik {(Crelles} Journal)},
	author = {Helmholtz, H.},
	year = {1858},
	pages = {25--55},
}

\bib{Kelvin1869}{article}{
	author={W. Thomson (Lord Kelvin)},
	title={On vortex motion},
	journal={Trans. Royal Soc. Edin.},
	number={25},
	pages={217--260},
	year={1869},
}

\bib{[Leandre-1999]}{article}{
   author={L{\'e}andre, R{\'e}mi},
   title={String structure over the Brownian bridge},
   journal={J. Math. Phys.},
   volume={40},
   date={1999},
   number={1},
   pages={454--479},
   issn={0022-2488},
   review={\MR{1657812 (2000k:58044)}},
   doi={10.1063/1.532781},
}

\bib{[Lyons-1998]}{article}{
   author={Lyons, Terry J.},
   title={Differential equations driven by rough signals},
   journal={Rev. Mat. Iberoamericana},
   volume={14},
   date={1998},
   number={2},
   pages={215--310},
   issn={0213-2230},
   review={\MR{1654527 (2000c:60089)}},
   doi={10.4171/RMI/240},
}

\bib{[LionsStFlour]}{book}{
   author={Lyons, Terry J.},
   author={Caruana, Michael},
   author={L{\'e}vy, Thierry},
   title={Differential equations driven by rough paths},
   series={Lecture Notes in Mathematics},
   volume={1908},
   note={Lectures from the 34th Summer School on Probability Theory held in
   Saint-Flour, July 6--24, 2004;
   With an introduction concerning the Summer School by Jean Picard},
   publisher={Springer},
   place={Berlin},
   date={2007},
   pages={xviii+109},
   isbn={978-3-540-71284-8},
   isbn={3-540-71284-4},
   review={\MR{2314753 (2009c:60156)}},
}
	
\bib{[LyonsQian-2002]}{book}{
   author={Lyons, Terry},
   author={Qian, Zhongmin},
   title={System control and rough paths},
   series={Oxford Mathematical Monographs},
   note={Oxford Science Publications},
   publisher={Oxford University Press},
   place={Oxford},
   date={2002},
   pages={x+216},
   isbn={0-19-850648-1},
   review={\MR{2036784 (2005f:93001)}},
   doi={10.1093/acprof:oso/9780198506485.001.0001},
}


\bib{[Rosenhead-1930]}{article}{
	title = {The Spread of Vorticity in the Wake Behind a Cylinder},
	volume = {127},
	issn = {1364-5021, 1471-2946},
	doi = {10.1098/rspa.1930.0078},
	number = {806},
	journal = {Proceedings of the Royal Society A: Mathematical, Physical and Engineering Sciences},
	author = {Rosenhead, L.},
	year = {1930},
	pages = {590--612},
}


\bib{[VincMeneguzzi-1991]}{article}{
	title = {The Spatial Structure and Statistical Properties of Homogeneous Turbulence},
	volume = {225},
	doi = {10.1017/S0022112091001957},
	journal = {Journal of Fluid Mechanics},
	author = {Vincent, A. and Meneguzzi, M.},
	year = {1991},
	pages = {1--20},
}


\bib{[Young-1936]}{article}{
   author={Young, L. C.},
   title={An inequality of the H\"older type, connected with Stieltjes
   integration},
   journal={Acta Math.},
   volume={67},
   date={1936},
   number={1},
   pages={251--282},
   issn={0001-5962},
   review={\MR{1555421}},
   doi={10.1007/BF02401743},
}

\end{biblist}
\end{bibdiv}

\end{document}
